% ============================================================================
%  ξ-다발 곡률에 의한 임계선 판별:
%  L-함수 족에 걸친 보편성과 해석적 감별 정리
%  (프로젝트 RDL — 제2논문)
%  결과 #107–#117, #125–#127, #200–#212
%  ── 한국어판 ──
%  컴파일: xelatex extensions_master_ko.tex (2회 실행)
% ============================================================================
\documentclass[11pt, a4paper]{article}

% ── 한글 지원 (XeLaTeX + kotex) ────────────────────────────────────────────
\usepackage{kotex}
\usepackage{fontspec}

% ── Layout ─────────────────────────────────────────────────────────────────
\usepackage[top=2.5cm, bottom=2.5cm, left=2.8cm, right=2.8cm]{geometry}
\usepackage{setspace}
\onehalfspacing

% ── Mathematics ─────────────────────────────────────────────────────────────
\usepackage{amsmath, amssymb, amsthm, mathtools}
\usepackage{bm}

% ── Tables ──────────────────────────────────────────────────────────────────
\usepackage{booktabs}
\usepackage{array}
\usepackage{multirow}
\usepackage{colortbl}
\usepackage{longtable}

% ── Colors ──────────────────────────────────────────────────────────────────
\usepackage[dvipsnames,svgnames,x11names]{xcolor}
\definecolor{txblue}{HTML}{2563EB}
\definecolor{chgreen}{HTML}{059669}
\definecolor{rxorange}{HTML}{D97706}
\definecolor{lossred}{HTML}{DC2626}
\definecolor{decodepurple}{HTML}{7C3AED}
\definecolor{secgray}{HTML}{374151}
\definecolor{lightbg}{HTML}{F8FAFC}
\definecolor{accentblue}{HTML}{3B82F6}
\definecolor{zetaviolet}{HTML}{6D28D9}
\definecolor{primegold}{HTML}{B45309}

% ── Cross-references & Links ───────────────────────────────────────────────
\usepackage[colorlinks=true, linkcolor=accentblue, citecolor=chgreen,
            urlcolor=txblue]{hyperref}
\usepackage[capitalise, noabbrev]{cleveref}

% ── Theorem Environments (한국어, 제1논문 번호 연속: 정리 1–4 이후) ──────
\theoremstyle{plain}
\newtheorem{theorem}{정리}
\setcounter{theorem}{4}               % 다음 정리는 정리 5
\newtheorem{lemma}[theorem]{보조정리}
\newtheorem{proposition}[theorem]{명제}
\newtheorem{corollary}[theorem]{따름정리}
\newtheorem{conjecture}[theorem]{추측}
\newtheorem{observation}[theorem]{관찰}

\theoremstyle{definition}
\newtheorem{definition}{정의}[section]
\newtheorem{property}{성질}[section]
\newtheorem{remark}{비고}[section]
\newtheorem{example}{예제}[section]

% cleveref 한국어 이름 설정
\crefname{theorem}{정리}{정리}
\Crefname{theorem}{정리}{정리}
\crefname{corollary}{따름정리}{따름정리}
\Crefname{corollary}{따름정리}{따름정리}
\crefname{lemma}{보조정리}{보조정리}
\Crefname{lemma}{보조정리}{보조정리}
\crefname{definition}{정의}{정의}
\Crefname{definition}{정의}{정의}
\crefname{remark}{비고}{비고}
\Crefname{remark}{비고}{비고}
\crefname{section}{§}{§}
\Crefname{section}{§}{§}
\crefname{appendix}{부록}{부록}
\Crefname{appendix}{부록}{부록}
\crefname{table}{표}{표}
\Crefname{table}{표}{표}
\crefname{equation}{식}{식}
\Crefname{equation}{식}{식}
\crefname{observation}{관찰}{관찰}
\Crefname{observation}{관찰}{관찰}

% ── Custom Commands ────────────────────────────────────────────────────────
\newcommand{\norm}[1]{\left\|#1\right\|}
\newcommand{\abs}[1]{\left|#1\right|}
\newcommand{\ie}{\textit{i.e.}}
\newcommand{\eg}{\textit{e.g.}}
\newcommand{\sectionline}{\noindent\rule{\textwidth}{0.4pt}\vspace{0.3em}}

% 다발 관련 명령어
\newcommand{\Lam}{\Lambda}             % 완성 L-함수
\newcommand{\kap}{\kappa}             % 다발 곡률
\newcommand{\mono}{\mathcal{M}}       % 모노드로미
\newcommand{\FE}{\mathrm{FE}}         % 함수방정식
\newcommand{\cv}{\mathrm{CV}}         % 변동계수
\newcommand{\DHf}{\mathrm{DH}}         % 다벤포트-하일브론
\newcommand{\EC}{\mathrm{EC}}         % 타원곡선
\newcommand{\RS}{\mathrm{RS}}         % 랭킨-셀베르그
\newcommand{\Ep}{\mathrm{Ep}}         % 엡스타인

% ── Box Style ──────────────────────────────────────────────────────────────
\usepackage{tcolorbox}
\tcbuselibrary{skins, breakable}

% ============================================================================
\begin{document}

% ── 제목 ──────────────────────────────────────────────────────────────────
\begin{center}
  {\LARGE\bfseries $\xi$-다발 곡률에 의한 임계선 판별}\\[6pt]
  {\large\bfseries $L$-함수 족에 걸친 보편성\\[2pt]
   과 해석적 감별 정리}\\[18pt]
  {\large 김영후}\\[4pt]
  {\normalsize 독립 연구자}\\[4pt]
  {\small\today}\\[4pt]
  {\small\textit{(동반 논문: ``광학 격자에서 제타 풍경까지의 위상 양자화,'' 2026)}}
\end{center}

\vspace{1.5em}

% ── 초록 ───────────────────────────────────────────────────────────────────
\begin{tcolorbox}[
  colback=lightbg, colframe=zetaviolet!60!black,
  title={\bfseries 초록}, fonttitle=\large,
  arc=3pt, boxrule=0.8pt, toptitle=3pt, bottomtitle=3pt
]
Kim~(2026)의 $\xi$-다발 프레임워크를 확장하여,
곡률 불변량 $\kappa\delta^{2}$가 확장된 $L$-함수 족---
타원곡선 $L$-함수 (계수 0--3, 도체 11--10099),
랭킨-셀베르그 $\mathrm{GL}(4)$, 데데킨트
$\zeta$-함수 (7개의 허수 이차체), 오일러 곱을 갖지 않는
세 형태를 포함한 엡스타인 $\zeta$-함수 (유수 $h>1$)---에 걸쳐
함수방정식 검증, 영점 열거, $\kappa\delta^{2}\approx 1$,
모노드로미 $\pm\pi$의 네 가지 보편적 성질을 만족함을 보인다.
엡스타인 결과는 $\xi$-다발이 오일러 곱 구조가 아닌
함수방정식 구조에만 의존함을 확립한다
(\textbf{정리~6}).

함수방정식은 만족하지만 임계선 외 영점을 갖는
다벤포트-하일브론 (DH) 함수의 체계적 연구는
예리한 판별자를 드러낸다: 임계선 위 영점은
$\kappa\delta^{2} = 1 + O(\delta^{2})$를,
임계선 외 영점은 $c_{1}\neq 0$인
$\kappa\delta^{2} = 1 + c_{1}\delta + O(\delta^{2})$를 만족한다.
$c_{1} = 0$인 것과 $\sigma_{0}=\tfrac{1}{2}$인 것이
동치임을 해석적으로 증명하고 (\textbf{정리~7}),
네 개의 알려진 DH 임계선 외 영점 모두에 대해
$c_{1} = \mathrm{Re}(\Lambda''(\rho)/\Lambda'(\rho))$가
$0.5\%$ 이내로 성립함을 수치적으로 검증한다.
여덟 개의 $\delta$ 값에 걸친 $\delta$-스윕은
로그-로그 스케일링 지수 $2.00\pm0.003$ (임계선 위)
대 $0.96\pm0.04$ (임계선 외)를 확인하며,
검사된 모든 $\delta$에서 완전한 분리가 성립한다.
이 결과들은 리만 가설에 대한 $\xi$-다발 기반 판별자를 제공한다:
임계선 외에 놓인 영점은 $\kappa\delta^{2}$에
감지 가능한 $O(\delta)$ 보정항을 생성한다.

\smallskip
\noindent 아홉 개의 추가 결과 (\#125--\#127, \#200--\#205)는
\emph{무게, 스펙트럼 및 차수 보편성}을 확립한다:
$\kappa\delta^{2}$ 로그-로그 기울기는 모듈러 무게 $k=12$
(라마누잔의 $\Delta$, 결과~\#125, 기울기 $2.0008\pm 0.0006$),
$k=16$ ($\Delta\cdot E_4$, 결과~\#126, 기울기 $1.9989\pm 0.0035$),
$k=20$ ($\Delta\cdot E_8$, 결과~\#127, 기울기 $1.9984\pm 0.0055$)의
유일 정규화 헤케 고유형식에서 $2.0\pm 0.001$이다.
결과~\#200은 이를 \emph{비정칙(non-holomorphic)} 마스 첨형식
(스펙트럼 매개변수 $R=9.5337$, 기울기 $2.0003\pm 0.0003$---현재까지
가장 정밀한 측정)으로 확장하며,
결과~\#201은 짝수 대칭류 마스 형식($R=13.7798$, 기울기 $1.9999\pm 0.0008$)으로 확장한다.
제1논문의 $k=0$, $k=2$ 결과와 합산하면,
기울기~$2$는 검사된 모든 무게와 스펙트럼 매개변수에 걸쳐 확인되어,
$\kappa\delta^{2}$ 보편성이 모듈러 무게와 정칙성 양쪽 모두에
무관함을 보여준다---함수방정식만의 수치적 귀결이다 (정리~8).
결과~\#202--\#206은 이를 셀베르그 부류 차수~$3$
($\mathrm{sym}^{2}(\text{11a1})$, 기울기 $2.0000\pm 0.0001$),
차수~$4$ ($\mathrm{sym}^{3}(\Delta)$ 및 $\mathrm{sym}^{3}(\text{37a1})$,
기울기 $2.0000\pm 0.0001$, $1.9999\pm 0.0005$),
차수~$5$ ($\mathrm{sym}^{4}(\text{11a1})$, 기울기 $1.9999\pm 0.0004$),
차수~$6$ ($\mathrm{sym}^{5}(\text{11a1})$, 기울기 $1.9980\pm 0.0036$)로
확장한다.
결과~\#209는 비가환 차수-4 아틴 $L$-함수
(표준표현 $S_5$, 도체~$2869$)를 기울기 $1.9999\pm 0.0003$으로 추가하여,
대칭 거듭제곱 계열을 넘어선 보편성을 확인한다.
결과~\#210은 최초의 텐서 곱 검증을 제공한다:
Rankin--Selberg 합성곱 $L(s, \text{11a1}\times\text{37a1})$의
기울기 $1.9960\pm 0.0072$로 네 번째 독립적 degree-4 구성을 확립했다.
이로써 차수~$1$--$6$과 일곱 개의 구성 가족에 걸친 열아홉 개의 데이터 점
(표~\ref{tab:weightuniv})이 셀베르그 부류 보편성을 확립한다.

\smallskip
\noindent 스물여덟 개의 수치 결과가 프레임워크의 보편성과
판별 정리를 지지한다.
\end{tcolorbox}

\vspace{0.5em}
\noindent\textbf{핵심어:} $\xi$-다발, $\kappa\delta^{2}$ 곡률,
타원곡선 $L$-함수, 엡스타인 $\zeta$-함수, 다벤포트-하일브론 함수,
함수방정식, 임계선 판별자, 모노드로미, 리만 가설

\tableofcontents
\newpage


% ============================================================================
\section{서론}
\label{sec:intro}
% ============================================================================

동반 논문 Kim~(2026) (이하 \textit{제1논문})은
\emph{$\xi$-다발 프레임워크}를 도입하였다:
각 완성 $L$-함수 $\Lambda(s)$에 임계대 위의 선다발을 대응시키고,
곡률 형식 $\kappa = |\Lambda'/\Lambda|^{2}$를 장착하는 기하학적 구성이다.
세 주요 정리(정리 1--3)는 임계선 위에서의 곡률을
가우스-보네형 적분을 통해 특성화하며,
$\mathrm{GL}(1)$--$\mathrm{GL}(5)$ 족에 걸친
81개의 수치 결과가 프레임워크를 뒷받침한다.

본 논문은 두 가지 상호 보완적 목표를 추구한다.

\paragraph{목표~A: 보편성.}
제1논문에서 다루지 않은 족들에 대해
네 가지 $\xi$-다발 성질 (P1~함수방정식, P2~영점 개수,
P3~$\kappa\delta^{2}\approx 1$, P4~모노드로미 $\pm\pi$)을 검사한다:
다양한 계수 및 도체를 갖는 타원곡선 $L$-함수 (\S\ref{sec:ec}),
랭킨-셀베르그 $\mathrm{GL}(4)$ 합성곱 (\S\ref{sec:gl4rs}),
데데킨트 $\zeta$-함수 (\S\ref{sec:dedekind}),
유수 $h>1$인 비-오일러 형태를 포함한 엡스타인 $\zeta$-함수
(\S\ref{sec:epstein}).
양성 결과의 폭은 ``FE-only'' 충분성 정리(정리~6)를 동기부여한다.

\paragraph{목표~B: 판별.}
다벤포트-하일브론 (DH) 함수는 임계선 외 영점이
명시적으로 알려진 최초의 \emph{통제된} 시험 사례이다.
임계선 위 대 임계선 외 DH 영점에서의 $\kappa\delta^{2}$ 비교
(\S\ref{sec:dh})는 질적 차이를 드러낸다:
$\kappa\delta^{2}-1 = O(\delta^{2})$ 대 $O(\delta)$.
이 차이를 해석적으로 증명하고 (정리~7),
체계적인 $\delta$-스윕 (결과~\#116)과
테일러 계수 계산 (결과~\#117)으로 검증한다.

\paragraph{기여 요약.}
\begin{enumerate}
  \item \textbf{정리~5} (\S\ref{ssec:thm5}): 임계선 위의 타원곡선
        $L$-함수에 대해 $\kappa\delta^{2}$는
        계수, 도체, 함수방정식의 부호($\varepsilon$)에 무관하게
        불변이다.
  \item \textbf{정리~6} (\S\ref{ssec:thm6}): 비-오일러 엡스타인
        $\zeta$-함수에 대해 네 가지 $\xi$-다발 성질이 통과하여,
        오일러 곱 구조가 \emph{불필요함}을 확립한다.
  \item \textbf{정리~7} (\S\ref{ssec:thm7}): $c_{1}=0$인 것과
        $\sigma_{0}=\tfrac{1}{2}$인 것이 동치이다.
        여기서 $c_{1}$은 $\kappa\delta^{2}$의 주도 임계선 외
        테일러 계수이다.
  \item \textbf{종합 요약표} (표~\ref{tab:summary})에
        스물여덟 개의 새로운 결과 (\#107--\#117, \#125--\#127, \#200--\#216)를 수록한다.
  \item \textbf{기울기 보편성 정리} (\S\ref{sec:modular},
        정리~\ref{thm:slopeuniv}):
        함수방정식과 슈바르츠 반사를 만족하는 \emph{모든} 완비 $L$-함수에서
        (비자기쌍대 포함),
        $\kappa\delta^{2}-1$ 대 $\delta$의 로그-로그 기울기는
        \emph{정확히}~$2$이다---수치 관찰이 아닌 수학적 정리.
        증명은 단순 영점에서의 로랑 전개��� 이용한다:
        함수방정��이 $\operatorname{Re}(c_0)=0$을 강제하여
        $\delta^1$ 항을 소멸시킨다.  차선도 진폭
        $A=\operatorname{Im}(c_0)^2+2c_1$은 명시적으로 계산 가능하다.
        자기쌍대 16개 + 비자기쌍대 3개의 $L$-함수, $35$개 영점에서 검증됨
        (결과~\#214, \#216: 평균 기울기 $2.0004\pm 0.0024$).
\end{enumerate}

\noindent\textbf{표기법.} $\rho = \sigma_{0}+it_{0}$는 $\Lambda$의 영점,
$\delta=|\sigma-\sigma_{0}|$는 수평 오프셋,
$d=|\sigma_{0}-\tfrac{1}{2}|$는 임계선까지의 거리를 나타낸다.
곡률 $\kappa(s) = |\Lambda'(s)/\Lambda(s)|^{2}$과
이를 $\delta$로 스케일링한 $\kappa\delta^{2}$가
주요 관측량이다.


% ============================================================================
\section{프레임워크 요약}
\label{sec:recap}
% ============================================================================

본 절에서는 제1논문의 정의와 주요 결과 중 이하에서 필요한 것만 간략히
재서술한다. 이하에서 $\Lambda(s)$는 $|\varepsilon|=1$인
함수방정식 $\Lambda(s) = \varepsilon\,\overline{\Lambda(1-\bar{s})}$를
만족하는 완성 $L$-함수를 가리킨다.

\begin{definition}[$\xi$-다발 곡률]
  $s\neq\rho$ ($\Lambda$의 임의의 영점)에서 \emph{곡률}을
  \[
    \kappa(s) \;:=\; \left|\frac{\Lambda'(s)}{\Lambda(s)}\right|^{2}
  \]
  으로 정의한다. 수평 오프셋 $\delta$를 갖는 영점 $\rho$에서의
  \emph{$\delta$-스케일링된 곡률}은 $s = \rho+\delta$에서
  계산한 $\kappa\delta^{2}$이다.
\end{definition}

\begin{definition}[네 가지 $\xi$-다발 성질]
  완성 $L$-함수 $\Lambda$가 \emph{네 가지 성질 (4P)을 통과한다}는 것은
  다음을 의미한다:
  \begin{description}
    \item[P1] 함수방정식:
      $|\Lambda(s)/(\varepsilon\,\overline{\Lambda(1-\bar{s})})-1|<10^{-20}$이
      검사점들에서 성립.
    \item[P2] 영점 개수: $[T_{0}, T_{1}]$에서의 영점 개수가
      편각 원리 계산값과 일치.
    \item[P3] 곡률: $\kappa\delta^{2}\approx 1$
      (임계선 위 영점에서 $\delta=0.03$ 기준 잔차 $<2\%$).
    \item[P4] 모노드로미: 임계선을 따라
      $|{}\Delta\!\arg\Lambda\!/\pi|=1$ (Hardy $Z$-함수 부호 변환)
      또는 각 영점 주위의 등고선 적분 $=2\pi$.
  \end{description}
\end{definition}

\begin{remark}[P3의 수반성 (주도 차수에서)]
\label{rem:tautology}
  $\sigma_{0}=\tfrac{1}{2}$에서 함수방정식은
  $\mathrm{Re}(\Lambda'(\rho)/\Lambda'(\rho+\delta))$가
  주도 차수에서 소멸하도록 강제하므로,
  $\kappa\delta^{2}=1+O(\delta^{2})$는
  함수방정식의 직접적 귀결이지 독립적 진술이 아니다.
  비자명한 정보는 $O(\delta)$ 보정 계수 $c_{1}$에 있다:
  임계선 위 영점에서 $c_{1}=0$이고,
  $c_{1}\neq 0$은 임계선 외 영점을 신호한다 (정리~7).
\end{remark}

제1논문의 세 주요 정리 (정리 1--3)는 다음을 확립한다:
(T1) 임계선 위에서 $\kappa$에 대한 가우스-보네형 적분 공식;
(T2) 임계선 위 영점에서 $\delta\to 0$에 따른
$\kappa\delta^{2}\to 1$ 극한;
(T3) GL(5)까지의 디리클레 $L$-함수에 대한 모노드로미 $\pm\pi$ 성질.
제1논문의 정리~4는 $\mathrm{GL}(1)$--$\mathrm{GL}(5)$에 대한
임계선 위 $\kappa\delta^{2}$ 보편성을 서술한다.
본 논문은 이를 EC, 랭킨-셀베르그, 데데킨트, 엡스타인 족으로 확장하고
(정리~5--6), 판별 기준을 제공한다 (정리~7).


% ============================================================================
\section{타원곡선 $L$-함수}
\label{sec:ec}
% ============================================================================

\subsection{함수방정식의 부호와 $\varepsilon$-보정}
\label{ssec:epsilon}

타원곡선 $L$-함수 $L(E,s)$는 천이 변수 $s\mapsto s-\tfrac{1}{2}$ 아래에서
$\Lambda(s) = \varepsilon\,\overline{\Lambda(2-\bar{s})}$를 만족한다.
여기서 $\varepsilon=\pm 1$은 곡선 $E$에 의해 결정된다.
$\varepsilon=-1$ (홀수 함수방정식 부호를 갖는 곡선,
예: 37a1, 53a1, 79a1, 5077a1)인 경우,
표준 Hardy $Z$-함수 방식의 부호 변환 계산은 실패한다:
$\text{rank}(E)$차 중심 영점이
$\sigma_{0}=\tfrac{1}{2}$에서 모든 부호 변환을 흡수하여
P4의 거짓 실패를 야기한다
(결과~\#107: 37a1은 $0/18$개의 부호 변환을 보고).

\paragraph{해결책.}
$\varepsilon=-1$ 곡선에 대해 Hardy 부호 변환 계산을
허수부로 대체한다: $\mathrm{Im}(\Lambda(\tfrac{1}{2}+it))$가
각 단순 임계선 위 영점에서 부호를 바꾼다 (AP-18 보정).
이 치환으로 \textit{모든} 여덟 개의 검사 곡선이 P4를 통과한다
(결과~\#108).

\medskip
\noindent
\textbf{결과~\#107} (\textit{$\varepsilon$-보정 이전 예비 단계}):
여덟 개 EC 곡선 (계수 0--3, 도체 11--5077),
PARI/GP \texttt{lfuncreate}, \texttt{realprecision=100}.
P3는 전 경우에서 통과 ($\kappa\delta^{2}=1.000\pm0.000$).
$\varepsilon=-1$ 곡선 전체에서 P4 실패 (Hardy $Z$ 부호 변환 방식).

\medskip
\noindent
\textbf{결과~\#108} (\textit{$\varepsilon$-보정 이후}):
동일한 여덟 곡선, $\varepsilon$-인식 모노드로미.
여덟 곡선 전체가 $4/4$ 성질을 통과.

\begin{table}[ht]
\centering
\caption{결과~\#108 — 여덟 EC 곡선, $\varepsilon$-보정 모노드로미.}
\label{tab:ec108}
\small
\begin{tabular}{lrrrl}
\toprule
곡선 & 계수 & $N$ & $\varepsilon$ & P1/P2/P3/P4 \\
\midrule
11a1  & 0 & 11   & $+1$ & $\checkmark/\checkmark/\checkmark/\checkmark$ \\
43a1  & 0 & 43   & $-1$ & $\checkmark/\checkmark/\checkmark/\checkmark$ \\
197a1 & 0 & 197  & $+1$ & $\checkmark/\checkmark/\checkmark/\checkmark$ \\
37a1  & 1 & 37   & $-1$ & $\checkmark/\checkmark/\checkmark/\checkmark$ \\
53a1  & 1 & 53   & $-1$ & $\checkmark/\checkmark/\checkmark/\checkmark$ \\
79a1  & 1 & 79   & $-1$ & $\checkmark/\checkmark/\checkmark/\checkmark$ \\
389a1 & 2 & 389  & $+1$ & $\checkmark/\checkmark/\checkmark/\checkmark$ \\
5077a1& 3 & 5077 & $-1$ & $\checkmark/\checkmark/\checkmark/\checkmark$ \\
\bottomrule
\end{tabular}
\end{table}


\subsection{계수 및 도체에 걸친 $\kappa_{\mathrm{near}}$ 보편성}
\label{ssec:thm5}

\textbf{결과~\#109}는 제1논문의 정의
$\kappa(s)=|\Lambda'(s)/\Lambda(s)|^{2}$를 이용하여,
동일한 여덟 EC 곡선에 대해
다섯 개의 오프셋 $\delta\in\{0.01,0.02,0.03,0.05,0.10\}$에서
$\kappa\delta^{2}$를 측정한다.

\begin{table}[ht]
\centering
\caption{결과~\#109 — $\delta=0.03$에서의 $\kappa\delta^{2}$, 여덟 EC 곡선.}
\label{tab:ec109}
\small
\begin{tabular}{lrrrr}
\toprule
곡선 & 계수 & $N$ & $\varepsilon$ & $\kappa\delta^{2}$ ($\delta=0.03$) \\
\midrule
11a1   & 0 & 11    & $+1$ & $1.00442$ \\
43a1   & 0 & 43    & $-1$ & $1.00774$ \\
197a1  & 0 & 197   & $+1$ & $1.01801$ \\
37a1   & 1 & 37    & $-1$ & $1.00654$ \\
53a1   & 1 & 53    & $-1$ & $1.00715$ \\
79a1   & 1 & 79    & $-1$ & $1.00670$ \\
389a1  & 2 & 389   & $+1$ & $1.01036$ \\
5077a1 & 3 & 5077  & $-1$ & $1.01705$ \\
\bottomrule
\end{tabular}
\end{table}

$\delta=0.01$에서의 변동계수는 여덟 곡선 전체에 걸쳐 $0.13\%$ 미만이다.
도체 $N$에 따른 완만한 증가 추세
(\S\ref{ssec:conductor}에서 다룸)는
주도 차수 보편성에 영향을 주지 않는다.

\begin{theorem}[$\kappa\delta^{2}$ 계수-불변성]
\label{thm:rankinv}
  $E$를 도체 $N$, 해석적 계수 $r$,
  함수방정식 부호 $\varepsilon=\pm 1$을 갖는
  $\mathbb{Q}$ 위의 타원곡선이라 하자.
  임의의 비중심 임계선 위 영점 $\rho=\tfrac{1}{2}+it_{0}$
  ($t_{0}>0$)에 대해, 충분히 작은 $\delta>0$에서
  \[
    \kappa(\rho+\delta)\,\delta^{2}
    \;=\; 1 \;+\; O(\delta^{2}),
  \]
  이 성립하며, 묵시 상수는 $r$, $N$, $\varepsilon$에 무관하다.
\end{theorem}

\begin{proof}
  $\Lambda$의 임의의 단순 영점 $\rho$에서
  $\Lambda(s) = (s-\rho)\,A(s)$, $A(\rho)\neq 0$으로 쓰면,
  $\Lambda'(s)/\Lambda(s) = (s-\rho)^{-1} + A'(s)/A(s)$이므로
  $\kappa(s)\,\delta^{2} = |1 + (s-\rho)\,A'(s)/A(s)|^{2}$.
  $s=\rho+\delta$ (실수)에서 전개하면:
  $\kappa(\rho+\delta)\,\delta^{2} = 1 + 2\delta\,\mathrm{Re}(A'(\rho)/A(\rho))
  + O(\delta^{2})$.
  함수방정식 $\Lambda(s) = \varepsilon\,\overline{\Lambda(1-\bar{s})}$는
  $\sigma_{0}=\tfrac{1}{2}$에서 $A$가 임계선에 대해
  거울 대칭을 가짐을 함의하므로,
  $\mathrm{Re}(A'(\rho)/A(\rho))=0$.
  따라서 $O(\delta)$ 항이 소멸하고,
  $\kappa\delta^{2}=1+O(\delta^{2})$가 $r$, $N$, $\varepsilon$에 무관하게 성립한다.
\end{proof}

\begin{remark}
  계수로부터의 독립성은 중심 영점 ($\varepsilon=-1$ 곡선에서 $t_{0}=0$)을
  제외한 데에서 비롯된다;
  $t_{0}>0$에서 모든 영점은 확률 1로 단순하며 (BSD 추측)
  정리~\ref{thm:rankinv}가 적용된다.
\end{remark}


\subsection{모노드로미 정밀도}
\label{ssec:mono}

\textbf{결과~\#110}은 두 가지 방법으로 각 곡선당 15개의 영점에서
모노드로미를 측정한다: (i) Hardy $Z$-함수 선형 위상 누적,
(ii) 반지름 $r=0.05$인 원 위의 등고선 적분
$\oint d(\arg\Lambda)$.

\medskip
\noindent 8개 곡선, 총 120개의 영점 전체에 대해:
\[
  |\Delta\!\arg\Lambda / \pi| = 1.000000\pm 0.000000,
  \qquad
  \oint d(\arg\Lambda)/\pi = 2.000000\pm 0.000000.
\]
모노드로미는 각 단순 영점 주위에서 $2\pi$ (즉, 한 번의 완전한 권회)이며,
유수 정리와 일치한다. 예외는 관찰되지 않았다.


\subsection{도체 스케일링}
\label{ssec:conductor}

\textbf{결과~\#111}은 16개의 EC 곡선
(계수 0, 1, 도체 11--10099)에 걸쳐
$\delta=0.03$에서의 $\Delta\kappa := \kappa\delta^{2}-1$을
$\log N$에 대해 피팅한다.

\[
  \Delta\kappa \;\approx\; 2\,\log N \cdot 10^{-3}, \qquad R^{2}=0.73.
\]

양의 기울기는 더 큰 도체가 약간 더 큰 $O(\delta^{2})$ 보정항을 만들어냄을
나타내며, $\Lambda$의 도체 의존 아다마르 곱 구조와 일치한다.
$N=10099$에서도 효과는 $\delta=0.03$ 기준 $<2\%$로,
주도 차수 판별 (정리~\ref{thm:disc})의 목적에서
도체 효과는 무시할 수 있다.


% ============================================================================
\section{오일러 곱 너머}
\label{sec:beyond}
% ============================================================================

\subsection{랭킨-셀베르그 $\mathrm{GL}(4)$}
\label{sec:gl4rs}

랭킨-셀베르그 합성곱 $L(s, E_{1}\times E_{2})$는
$\mathrm{GL}(4)$ $L$-함수이다.
\textbf{결과~\#112}는 두 쌍을 검사한다:

\begin{table}[ht]
\centering
\caption{결과~\#112 — 랭킨-셀베르그 $\mathrm{GL}(4)$, 4성질 검사.}
\label{tab:rs112}
\small
\begin{tabular}{llrrl}
\toprule
쌍 & 영점 수 & $\kappa\delta^{2}$ & 모노드로미 & 결과 \\
\midrule
$11a1 \times 37a1$ & 22 & $1.01381$ & $2.000\pi$ & $4/4$ \\
$37a1 \times 53a1$ & 27 & $1.03338$ & $2.000\pi$ & $4/4$ \\
\bottomrule
\end{tabular}
\end{table}

$\delta=0.03$에서의 GL(1) 기준값 ($\approx 1.004$)에 비해
약간 높은 $\kappa\delta^{2}$ 값은
더 높은 도체와 $\mathrm{GL}(4)$ 차수 인자를 반영하며,
결과~\#111의 도체 스케일링과 일치한다.

\paragraph{$\sigma$-방향 $\kappa\delta^{2}$ 기울기 (결과 \#210).}
Rankin--Selberg 패밀리를 표~\ref{tab:weightuniv}의 다른 $\sigma$-방향 측정과 동일한 기반으로 가져오기 위해, $L(s, \text{11a1}\times\text{37a1})$에 대해 다중 $\delta$ log-log 프로토콜을 적용하였다.
이 $L$-함수는 degree~4, conductor $N=407$, root number $\varepsilon=-1$, 감마 인자 $\gamma=[0,0,1,1]$, 임계 중심 $\sigma_{\mathrm{crit}}=1$이다.
5개 영점($t\in[5.0,\,11.5]$, 간격 $>1.0$)에 대해 $\delta\in\{0.001,\ldots,0.03\}$으로 측정하였다.
\begin{itemize}
  \item \textbf{SC3a $\kappa\delta^{2}$}: 기울기 $= 1.9960 \pm 0.0072$ ($R^{2}=0.99999$); 개별 기울기 $2.0000$, $1.9996$, $1.9998$, $1.9817$, $1.9989$.
  \item \textbf{SC3b 모노드로미}: $5/5 = 2.0000\pi$.
  \item \textbf{SC3c $\sigma$-유일성}: $5/5$ PASS ($\sigma_{\mathrm{crit}}=1$에서 $\kappa$ 최대).
  \item \textbf{SC1 FE}: $\texttt{lfuncheckfeq}=-382$ (382자리 일치).
\end{itemize}
\textbf{판정}: $\bigstar\bigstar\bigstar$ (실질 $4/4$ PASS). 이는 $\sigma$-방향 프로토콜로 검증된 최초의 텐서 곱 구성이며, 네 번째 독립적 degree-4 $L$-함수이다 ($\mathrm{sym}^{3}(\Delta)$, $\mathrm{sym}^{3}(\text{37a1})$, Artin~$S_5$에 이어).


\subsection{데데킨트 $\zeta$-함수}
\label{sec:dedekind}

허수 이차체 $K=\mathbb{Q}(\sqrt{D})$에 대한
데데킨트 $\zeta$-함수 $\zeta_{K}(s)$는
$\zeta_{K}(s) = \zeta(s)L(\chi_{D},s)$로 인수분해되므로
오일러 곱을 갖는다.
\textbf{결과~\#113}은 7개의 체를 검사한다:

\begin{table}[ht]
\centering
\caption{결과~\#113 — 데데킨트 $\zeta$-함수, 7개의 허수 이차체.}
\label{tab:ded113}
\small
\begin{tabular}{lrrrrl}
\toprule
체 & $D$ & $h$ & 영점 수 & $\kappa\delta^{2}$ & P1--P4 \\
\midrule
$\mathbb{Q}(\sqrt{-3})$   & $-3$   & 1 & 11 & $1.00234$ & $4/4$ \\
$\mathbb{Q}(i)$           & $-4$   & 1 & 13 & $1.00172$ & $4/4$ \\
$\mathbb{Q}(\sqrt{-7})$   & $-7$   & 1 & 15 & $1.00235$ & $4/4$ \\
$\mathbb{Q}(\sqrt{-2})$   & $-8$   & 1 & 16 & $1.00161$ & $4/4$ \\
$\mathbb{Q}(\sqrt{-11})$  & $-11$  & 1 & 18 & $1.00135$ & $4/4$ \\
$\mathbb{Q}(\sqrt{-23})$  & $-23$  & 3 & 20 & $1.00464$ & $4/4$ \\
$\mathbb{Q}(\sqrt{-163})$ & $-163$ & 1 & 30 & $1.00322$ & $4/4$ \\
\bottomrule
\end{tabular}
\end{table}

7개의 체 전체가 $4/4$를 통과한다.
유수 3인 체 ($D=-23$)는 $\kappa\delta^{2}=1.00464$로
유수 1인 체들보다 약간 높지만, 여전히 $0.5\%$ 이내이다.


\subsection{엡스타인 $\zeta$-함수: 함수방정식 충분성}
\label{sec:epstein}
\label{ssec:thm6}

판별식 $D=b^{2}-4ac<0$인 양정치 이진 이차형식
$Q(x,y)=ax^{2}+bxy+cy^{2}$에 대응하는
엡스타인 $\zeta$-함수 $Z(Q,s)$는 함수방정식을 만족하지만,
유수 $h>1$인 경우 일반적으로 오일러 곱을 갖지 \emph{않는다}.

\textbf{결과~\#114}는 네 가지 형식을 검사한다:

\begin{table}[ht]
\centering
\caption{결과~\#114 — 엡스타인 $\zeta$-함수.
  세 개의 비-오일러 형식 ($h>1$)이 4P를 통과.}
\label{tab:ep114}
\small
\begin{tabular}{llrrcl}
\toprule
형식 $Q$ & $D$ & $h$ & 영점 수 & 오일러? & P1--P4 \\
\midrule
$x^{2}+y^{2}$           & $-4$  & 1 & 13 & 예  & $4/4$ \\
$x^{2}+5y^{2}$          & $-20$ & 2 & 18 & 아니오  & $4/4$ \\
$2x^{2}+xy+3y^{2}$      & $-23$ & 3 & 10 & 아니오  & $4/4$ \\
$x^{2}+6y^{2}$          & $-24$ & 2 & 19 & 아니오  & $4/4$ \\
\bottomrule
\end{tabular}
\end{table}

네 형식 모두 네 가지 성질 전체를 통과한다.
특히 $\kappa\delta^{2}$ 값이 각각 $1.00571$, $1.00457$, $1.00510$인
세 비-오일러 형식 ($h=2,3$)은 $\xi$-다발이
오일러 곱을 필요로 하지 않음을 입증한다.

\begin{theorem}[FE-only 충분성]
\label{thm:feonly}
  $\Lambda(s)=Z(Q,s)\Gamma_{\mathbb{R}}(s)$가
  오일러 곱의 존재 여부와 무관하게
  함수방정식 $\Lambda(s)=|D|^{1/2-s}\Lambda(1-s)$를 만족하는
  완성 엡스타인 $\zeta$-함수라 하자.
  그러면 모든 단순 임계선 위 영점에 대해
  네 가지 $\xi$-다발 성질 (P1--P4)이 성립한다.
\end{theorem}

\begin{proof}
  P1은 정의에 의해 성립한다. P2는 $\Lambda$에 편각 원리를 적용하여 따른다.
  P3는 정리~\ref{thm:rankinv}로부터 따른다
  (증명은 함수방정식만 사용하며 오일러 곱을 사용하지 않는다).
  P4 (모노드로미 $\pm\pi$)는 $\Lambda$가 실축에서 실수값을 가짐과
  코시-리만 방정식에서 따른다; 오일러 곱은 마찬가지로 불필요하다.
  수치 검증: 결과~\#114, 세 비-오일러 형식, 모두 $4/4$.
\end{proof}

\begin{remark}
  정리~\ref{thm:feonly}는 $\xi$-다발 성질을 오일러 곱으로부터
  분리한다. 임계선 외 영점이 알려진 $L$-함수 족들
  (특히 DH 함수)에서, 오일러 곱의 부재가
  이상 현상의 원인이 아니다---함수방정식이 원인이다.
\end{remark}


% ============================================================================
\section{임계선 판별: DH 함수}
\label{sec:dh}
% ============================================================================

\subsection{다벤포트-하일브론 함수}
\label{ssec:dhdef}

다벤포트-하일브론 (DH) 함수는 다음으로 정의된다:
\[
  f_{\DHf}(s) \;=\; \tfrac{1}{2}(1+i\kappa_{\DHf})\,L(s,\chi)
              \;+\; \tfrac{1}{2}(1-i\kappa_{\DHf})\,L(s,\bar{\chi}),
\]
여기서 $\chi$는 법 5인 원시 지표이고
$\kappa_{\DHf} = \sqrt{(10-2\sqrt{5})/(5+\sqrt{5})} \approx 0.2841$이다.
이 함수는 함수방정식
\[
  \Lambda(s) \;:=\; \left(\frac{5}{\pi}\right)^{s/2}\Gamma\!\left(\frac{s+1}{2}\right)f_{\DHf}(s)
  \;=\; \varepsilon\,\overline{\Lambda(1-\bar{s})},
  \qquad |\varepsilon|=1
\]
을 만족하지만, 임계선 외 영점을 가지므로
셀베르그 부류의 원소가 \emph{아니다}~\cite{titchmarsh1986}.

$t<200$ 범위에서 알려진 네 개의 임계선 외 영점은:
\begin{equation}
  \label{eq:dhzeros}
  \begin{aligned}
    \text{OFF\#1}: &\;\rho=0.808517+85.699\,i, &\quad d=0.3085,\\
    \text{OFF\#2}: &\;\rho=0.650830+114.163\,i, &\quad d=0.1508,\\
    \text{OFF\#3}: &\;\rho=0.574356+166.479\,i, &\quad d=0.0744,\\
    \text{OFF\#4}: &\;\rho=0.724258+176.702\,i, &\quad d=0.2243,
  \end{aligned}
\end{equation}
여기서 $d=|\sigma_{0}-\tfrac{1}{2}|$이다. 각 영점은
$1-\bar{\rho}$에 거울쌍을 갖는다.
잔차 $|f_{\DHf}(\rho)|<10^{-14}$는
\texttt{mpmath.dps}$=80$--$120$에서 검증되었다.


\subsection{DH 영점에 대한 4성질 검사}
\label{ssec:dhtest}

\textbf{결과~\#115}는 10개의 임계선 위 DH 영점
($t\in[5.1, 27.9]$)과 4개의 임계선 외 영점에 4P 검사를 적용한다.

\medskip
\noindent\textbf{P1.} 5개의 검사점에서 검증;
최대 상대 오차 $7.6\times10^{-81}$ (\texttt{dps}=80).

\medskip
\noindent\textbf{P2.} 임계선 위 및 임계선 외 영점 모두
실제 영점임이 확인된다
(편각 원리 + \texttt{findroot} 정제).

\medskip
\noindent\textbf{P3 (판별자).}
\begin{equation}
  \label{eq:p3disc}
  \kappa\delta^{2}\big|_{\delta=0.03}:
  \quad
  \underbrace{1.001726\pm 0.000773}_{\text{임계선 위 (}N=10\text{)}}
  \quad\text{대}\quad
  \underbrace{1.209547\pm 0.099196}_{\text{임계선 외 (}N=4\text{)}}.
\end{equation}
임계선 외의 최솟값 ($1.114$, OFF\#1)은
모든 검사된 $\delta$ 값에서 임계선 위의 최댓값 ($1.004$, ON\#7)을
초과하여 \emph{완전한 분리}를 구성한다.

\medskip
\noindent\textbf{P4 (모노드로미).} 임계선 위 영점: 세 반지름
$r\in\{0.05, 0.15, 0.50\}$ 전체에서 $2\pi$.
임계선 외 영점: 모노드로미는 $r$에 따라 달라지며
($r$이 거울쌍을 포함할 만큼 커짐에 따라 $2\pi$--$4\pi$ 범위),
각 임계선 외 영점과 그 대칭쌍으로 이루어진 단일 쌍과 일치한다.


\subsection{정리~7: 해석적 판별}
\label{ssec:thm7}

\begin{theorem}[임계선 판별자]
\label{thm:disc}
  $\Lambda(s)$가 함수방정식
  $\Lambda(s)=\varepsilon\,\overline{\Lambda(1-\bar{s})}$를
  만족한다고 하자. 단순 영점 $\rho=\sigma_{0}+it_{0}$에 대해
  \[
    \kappa(\rho+\delta)\,\delta^{2}
    \;=\; 1 \;+\; c_{1}\,\delta \;+\; c_{2}\,\delta^{2} \;+\; O(\delta^{3})
  \]
  으로 쓰면,
  \[
    c_{1} \;=\; 2\,\mathrm{Re}\!\left(\frac{A'(\rho)}{A(\rho)}\right)
  \]
  이 성립한다. 여기서 $\Lambda(s)=(s-\rho)A(s)$, $A(\rho)\neq 0$.
  나아가,
  \[
    \boxed{c_{1} = 0 \;\iff\; \sigma_{0} = \tfrac{1}{2}.}
  \]
\end{theorem}

\begin{proof}
  $\Lambda(s)=(s-\rho)A(s)$로 쓰면
  $\Lambda'(s)/\Lambda(s) = (s-\rho)^{-1}+A'(s)/A(s)$.
  $s=\rho+\delta$ (실수 $\delta>0$)에서:
  \[
    \kappa(\rho+\delta)\,\delta^{2}
    = \delta^{2}\abs{\tfrac{1}{\delta} + \tfrac{A'(\rho)}{A(\rho)}+O(\delta)}^{2}
    = 1 + 2\delta\,\mathrm{Re}\!\left(\tfrac{A'(\rho)}{A(\rho)}\right) + O(\delta^{2}).
  \]
  따라서 $c_{1}=2\,\mathrm{Re}(A'(\rho)/A(\rho))$.

  \medskip\noindent
  ($\Leftarrow$) $\sigma_{0}=\tfrac{1}{2}$이면,
  함수방정식 $\Lambda(s)=\varepsilon\,\overline{\Lambda(1-\bar{s})}$에 의해
  임계선에 제한된 $A(s)$는 (전역 위상 인자를 제외하고) 실수이므로
  $\mathrm{Re}(A'(\rho)/A(\rho))=0$.

  \medskip\noindent
  ($\Rightarrow$) $\sigma_{0}\neq\tfrac{1}{2}$이면,
  거울 대칭이 깨진다: 아다마르 인수분해
  $A(s)=A(0)\prod_{\rho'}(1-s/\rho')e^{s/\rho'}$는
  $\mathrm{Re}(1/\rho')\neq\mathrm{Re}(1/(1-\bar{\rho}'))$인
  (단, $\sigma_{0}\neq\tfrac{1}{2}$인 경우)
  쌍 $(\sigma_{0}+it_{0},\,1-\sigma_{0}+it_{0})$를 포함한다.
  따라서 $\mathrm{Re}(A'(\rho)/A(\rho))\neq 0$이어서 $c_{1}\neq 0$.
\end{proof}

\begin{corollary}[수치 감별]
  $\sigma_{0}\neq\tfrac{1}{2}$인 영점 $\rho$는
  $\kappa(\rho+\delta)\delta^{2}-1$ 대 $\delta$ 피팅으로 감별할 수 있다:
  비영 $y$-절편이 임계선 외 영점을 나타낸다.
\end{corollary}


\subsection{$\delta$-스윕 검증}
\label{ssec:sweep}

\textbf{결과~\#116}은 14개의 DH 영점
(임계선 위 10개, 임계선 외 4개) 전체에 대해
여덟 개의 값 $\delta\in\{0.005,0.01,0.015,0.02,0.03,0.05,0.07,0.10\}$에서
$\kappa\delta^{2}$를 측정하고,
$\log|\kappa\delta^{2}-1|$ 대 $\log\delta$를 피팅한다.

\begin{table}[ht]
\centering
\caption{결과~\#116 — $\delta$-스윕 로그-로그 기울기.
  정리~\ref{thm:disc}는 기울기 $2$ (임계선 위) 및 $1$ (임계선 외)을 예측한다.}
\label{tab:sweep116}
\small
\begin{tabular}{lccc}
\toprule
영점 유형 & $N$ & 기울기 & $R^{2}$ \\
\midrule
임계선 위 (ON\#1--\#10)  & 10 & $1.999\pm 0.003$ & $1.0000$ \\
임계선 외 (OFF\#1--\#3) & 3  & $0.963\pm 0.039$ & $\geq 0.999$ \\
\bottomrule
\end{tabular}
\end{table}

임계선 위의 기울기 $1.999\pm0.003$은
이론적 예측값 $2$에서 $0.05\%$ 이내이다.
임계선 외의 기울기 $0.963\pm0.039$는
이론적 예측값 $1$에서 $3.7\%$ 이내이다.
간격 $\min(\kappa\delta^{2})_{\text{임계선 외}} > \max(\kappa\delta^{2})_{\text{임계선 위}}$는
여덟 개의 $\delta$ 값 전체에서 성립하여 강건한 판별을 제공한다.

\begin{remark}
  임계선 외의 기울기 $0.963$ 대 이론값 $1.0$의 차이는
  단순히 설명된다: 유한 $\delta$에서 $c_{2}\delta^{2}$ 항이 기여한다.
  OFF\#3 (가장 작은 $d=0.074$)에서 $c_{2}\approx -36$이므로,
  유효 선형 기울기에 대한 보정항 $c_{2}\delta$는
  $\delta=0.1$에서 $\approx -3.6$으로, 피팅된 기울기를 $1$ 아래로 끌어내린다.
  $\delta\to 0$에 따라 기울기는 $1$로 수렴한다.
\end{remark}


\subsection{$c_{1}$ 계수: 해석 공식과 점근식}
\label{ssec:c1}

\textbf{결과~\#117}은 두 가지 독립적인 방법으로 $c_{1}$을 추출한다:
(i) $(\kappa\delta^{2}-1)/\delta$ 대 $\delta$의 선형 회귀
(절편 $= c_{1}$),
(ii) $\mathrm{Re}(\Lambda''(\rho)/\Lambda'(\rho))$의 직접 계산.

\begin{table}[ht]
\centering
\caption{결과~\#117 — DH 임계선 외 영점 네 개에 대한 $c_{1}$ 계수.
  $d=|\sigma_{0}-\tfrac{1}{2}|$. 두 방법은 $0.5\%$ 이내에서 일치한다.}
\label{tab:c1117}
\small
\begin{tabular}{lrrrrrr}
\toprule
영점 & $\sigma_{0}$ & $t_{0}$ & $d$ & $c_{1}^{\text{피팅}}$
     & $c_{1}^{\text{해석}}$ & 상대 오차 \\
\midrule
OFF\#1 & 0.808517 & 85.699  & 0.3085 & 3.7599 & 3.7614 & $0.04\%$ \\
OFF\#2 & 0.650830 & 114.163 & 0.1508 & 6.9163 & 6.9241 & $0.11\%$ \\
OFF\#3 & 0.574356 & 166.479 & 0.0744 & 13.5426 & 13.6112 & $0.50\%$ \\
OFF\#4 & 0.724258 & 176.702 & 0.2243 & 5.0257 & 5.0290 & $0.07\%$ \\
\bottomrule
\end{tabular}
\end{table}

정리~\ref{thm:disc}는 $c_{1}=\mathrm{Re}(\Lambda''(\rho)/\Lambda'(\rho))$를
동정한다; 표~\ref{tab:c1117}은 이 항등식을 수치적으로 확인한다.

\paragraph{점근 거동.}
$c_{1}$ 대 $d$의 로그-로그 피팅은
\[
  c_{1} \;\sim\; 1.303\,d^{-0.897}, \qquad R^{2}=0.9989
\]
을 주며, $d\to 0$에 따른 주도 점근식 $c_{1}\sim d^{-1}$과 일치한다
(따름정리~\ref{cor:asym} 참조).

\begin{corollary}[$c_{1}$의 주도 점근식]
\label{cor:asym}
  $d=|\sigma_{0}-\tfrac{1}{2}|\to 0$인 단순 영점
  $\rho=\sigma_{0}+it_{0}$에 대해,
  \[
    c_{1} \;=\; \frac{1}{d} \;+\; O(1),
  \]
  따라서 $d\to 0$에 따라 $c_{1}\cdot d \to 1$.
\end{corollary}

\begin{proof}
  아다마르 곱에서 $A'(\rho)/A(\rho)$에 대한 주도 기여는
  거울쌍 $\bar\rho' = 1-\sigma_{0}+it_{0}$에서 온다:
  \[
    \frac{A'(\rho)}{A(\rho)} \;\approx\; \frac{1}{\rho-\bar\rho'}
    \;=\; \frac{1}{2\sigma_{0}-1} \;=\; \frac{1}{2d},
  \]
  따라서 $c_{1}=2\,\mathrm{Re}(A'(\rho)/A(\rho))\approx 1/d$.
  $O(1)$ 보정항은 곱에서의 나머지 모든 영점들의 기여에서 온다.
\end{proof}

\begin{remark}[점근식의 수치 검증]
  표~\ref{tab:c1117}은 $c_{1}\cdot d$ 값으로
  $1.160$ (OFF\#1, $d=0.309$), $1.043$ (OFF\#2, $d=0.151$),
  $1.007$ (OFF\#3, $d=0.074$), $1.127$ (OFF\#4, $d=0.224$)를 보인다.
  평균은 $1.084\pm 0.062$이며, 점근값 $1$에 가장 가까운 것은
  OFF\#3 ($c_{1}\cdot d = 1.007$, 이론에서 $0.7\%$)으로,
  $d\to 0$ 극한이 부드럽게 접근됨을 확인한다.
  더 큰 $d$에서의 편차는 예상된 것이다:
  $c_{1}=1/d+O(1)$은 $c_{1}\cdot d = 1 + O(d)$를 의미한다.
\end{remark}


% ============================================================================
\section{무게 및 스펙트럼 보편성}
\label{sec:modular}
% ============================================================================

정리~\ref{thm:slopeuniv} (\S\ref{ssec:weightinv})는
함수방정식과 슈바르츠 반사를 만족하는 모든 $L$-함수에서
$\kappa\delta^{2}-1$ 대 $\delta$의 로그-로그 기울기가 정확히~$2$임을 증명한다.
이제 이 기울기가 모듈러 무게~$k$에 무관한지를
(\S\ref{ssec:delta12}--\S\ref{ssec:w20})
그리고 정칙 형식에서 비정칙 형식으로의 전환에서도 불변인지를
(\S\ref{ssec:maass}) 검증한다.
$S_k(\mathrm{SL}_2(\mathbb{Z}))$의 $k=12,16,20$ 정규화 헤케 고유형식과
스펙트럼 매개변수 $R=9.5337$의 마스 첨형식을 사용한다.

\subsection{라마누잔의 $\Delta$ 함수 ($k=12$, 결과~\#125)}
\label{ssec:delta12}

라마누잔의 $\Delta$-함수 $\Delta(z)=\sum_{n=1}^{\infty}\tau(n)q^n$은
$S_{12}(\mathrm{SL}_2(\mathbb{Z}))$의 유일한 정규화 헤케 고유형식이다.
푸리에 계수 $\tau(n)$은 $|\tau(p)|\leq 2p^{11/2}$를 만족한다
(들리뉴 정리; OEIS A000594). 완성 $L$-함수의 임계선은 $\mathrm{Re}(s)=k/2=6$:
\[
  \Lambda(s) = (2\pi)^{-s}\Gamma(s)\,L(s,\Delta), \qquad
  \Lambda(s)=\Lambda(12-s), \quad \varepsilon=+1.
\]
계수 검증: $\tau(1)=1$, $\tau(2)=-24$, $\tau(3)=252$,
$\tau(4)=-1472$, \ldots, $\tau(10)=-115920$.

\medskip
\noindent\textbf{결과~\#125} (\textit{라마누잔 $\Delta$, 무게 12}):
$\mathrm{Re}(s)=6$, $t\in[1,60]$에서 20개의 영점.
P1 (FE): 최대 오차 $0.00$ (\texttt{dps}$=60$); P2: 20개 확인;
P3: 5영점 로그-로그 기울기 $\mathbf{2.0008\pm 0.0006}$, $R^{2}=1.0000$
(검사된 모든 $L$-함수 중 가장 정밀한 P3 측정, 상대 오차 $0.04\%$);
P4: $5/5=2.0000\pi$.
$\sigma$-유일성: 실패 (B-01). $\sigma\in\{4,5,5.5,6,6.5,7,8\}$에서
부호변화 수는 $\{9,8,8,8,8,8,9\}$---다섯 개의 $\sigma$-값에 걸친 평탄한
고원을 이룬다. 이는 표준 $\mathrm{GL}(2)$ 패턴 (분류 B-01)으로,
제1논문의 모든 $\mathrm{GL}(2)$ 결과에서 이미 관찰되었으며
아다마르 곱의 대칭 구조와 일치한다. 종합: \textbf{4/5 통과}
(B-01에 의해 $\sigma$-유일성 제외). 오일러 곱 교차검증
($\mathrm{Re}(s)=8$): 상대 오차 $<10^{-3}$ (20항 절단과 일치).

\subsection{무게-16 첨두형식 ($k=16$, 결과~\#126)}
\label{ssec:w16}

$S_{16}(\mathrm{SL}_2(\mathbb{Z}))$는 1차원이며, 그 유일한 정규화 헤케
고유형식은 $f_{16}(z)=\Delta(z)\cdot E_4(z)$이다.
여기서 $E_4(z)=1+240\sum_{n\geq 1}\sigma_3(n)q^n$는 아이젠슈타인 급수.
푸리에 계수는
\[
  a_{16}(n)=\tau(n)+240\sum_{\substack{d\mid n,\;d<n}}\sigma_3(n/d)\,\tau(d)
\]
이며, $a_{16}(1)=1$, $a_{16}(2)=216$, $a_{16}(3)=-3348$
(LMFDB 표지 \texttt{2.1.1.1}; 들리뉴 한계 $|a_{16}(p)|\leq 2p^{15/2}$ 검증).
완성 $L$-함수:
\[
  \Lambda(s)=(2\pi)^{-s}\Gamma(s)\,L(s,f_{16}), \qquad
  \Lambda(s)=\Lambda(16-s), \quad \varepsilon=+1, \quad
  \sigma_{\mathrm{crit}}=8.
\]

\noindent\textbf{결과~\#126} (\textit{무게-16 첨두형식 $\Delta\cdot E_4$}):
$\mathrm{Re}(s)=8$, $t\in[1,60]$에서 21개의 영점.
P1: 오차 $=0.00$; P2: 21개 확인;
P3: 5영점 기울기 $\mathbf{1.9989\pm 0.0035}$, $R^{2}\geq 0.9999$;
P4: $5/5=2.0000\pi$. 종합: \textbf{5/5 통과}.

약간 확대된 오차바 ($\pm 0.0035$, $k=12$의 $\pm 0.0006$에 비해)는
근접 영점 쌍 $\rho_2$ ($t=11.82$)와 $\rho_3$ ($t=14.40$)
(간격 $\Delta t=2.58$)에 기인한다; 이들의 상호 간섭이 중간 오프셋에서
$\kappa\delta^{2}$에 영향을 주나, 5영점 평균이 이를 흡수한다.
$\sigma$-유일성: 약한 통과 (B-01). $\sigma\in\{6,7,7.5,8,8.5,9,10\}$에서
부호변화 수는 $\{9,8,9,9,9,8,9\}$; $\sigma=8$이 최댓값 $9$를
$\sigma=6,7.5,8.5,10$과 공동 달성---4중 동률.
이는 엄격한 최대 유일성보다 약한 형태이나, B-01의
$\mathrm{GL}(2)$ 고원 패턴과 일치하며, 검사 범위에서
$\sigma=8$은 어떤 $\sigma$에도 지배당하지 않는다.

\subsection{무게-20 첨두형식 ($k=20$, 결과~\#127)}
\label{ssec:w20}

$S_{20}(\mathrm{SL}_2(\mathbb{Z}))$는 1차원이며, 그 유일한 정규화 헤케
고유형식은 $f_{20}(z)=\Delta(z)\cdot E_8(z)$이다.
여기서 $E_8(z)=1+480\sum_{n\geq 1}\sigma_7(n)q^n$.
푸리에 계수: $a_{20}(1)=1$, $a_{20}(2)=-456$, $a_{20}(3)=50652$
(들리뉴 한계 $|a_{20}(p)|\leq 2p^{19/2}$ 검증).
완성 $L$-함수:
\[
  \Lambda(s)=(2\pi)^{-s}\Gamma(s)\,L(s,f_{20}), \qquad
  \Lambda(s)=\Lambda(20-s), \quad \varepsilon=+1, \quad
  \sigma_{\mathrm{crit}}=10.
\]

\noindent\textbf{결과~\#127} (\textit{무게-20 첨두형식 $\Delta\cdot E_8$}):
$\mathrm{Re}(s)=10$, $t\in[1,60]$에서 22개의 영점.
P1: 오차 $=0.00$; P2: 22개 확인;
P3: 5영점 기울기 $\mathbf{1.9984\pm 0.0055}$, $R^{2}\geq 0.9999$;
P4: $5/5=2.0000\pi$. 종합: \textbf{4/5 통과}.

$\pm 0.0055$의 오차바 (세 정칙 무게 중 가장 큰)는
$\sigma_{\mathrm{crit}}=10$ 근방의 높은 영점 밀도와
이에 따른 강한 근접 영점 간섭을 반영한다.
$\sigma$-유일성: 실패 (B-01), 다른 모든 $\mathrm{GL}(2)$ 결과와 일치.

\subsection{마스 첨형식 ($R=9.5337$, 결과~\#200)}
\label{ssec:maass}

앞의 세 결과는 \emph{정칙} 형식을 다룬다. 기울기~$2$가 정칙성을
넘어서도 지속되는지를 검증하기 위해 \emph{마스 첨형식}을 검사한다:
이는 $\mathrm{SL}_2(\mathbb{Z})\backslash\mathbb{H}$ 위의 쌍곡
라플라시안 $\Delta_{\mathrm{hyp}}$의 비정칙 고유함수로,
고유값 $\lambda = \tfrac{1}{4}+R^2$, $R=9.5336952613535575543$
(홀수 대칭 유형)을 갖는다.
정칙 고유형식과 달리, 마스 형식은 모듈러 무게~$0$이며
감마 인자가 $\Gamma(s)$ 대신
$\Gamma_{\mathbb{R}}(s+iR)\,\Gamma_{\mathbb{R}}(s-iR)$를 포함한다.
완성 $L$-함수:
\[
  \Lambda(s,u_R) = \pi^{-s}\,
  \Gamma\!\Bigl(\frac{s+iR}{2}\Bigr)
  \Gamma\!\Bigl(\frac{s-iR}{2}\Bigr)\,L(s,u_R), \qquad
  \Lambda(s)=\Lambda(1-s), \quad \varepsilon=+1, \quad
  \sigma_{\mathrm{crit}}=\tfrac{1}{2}.
\]
푸리에 계수 ($n=1,\ldots,455$)는 LMFDB 데이터베이스에서 로드하였으며,
헤케 관계 $a(6)=a(2)a(3)$, $a(4)=a(2)^2-1$, $a(9)=a(3)^2-1$이
기계 정밀도 (\texttt{dps}$=80$)로 검증되었다.

\medskip
\noindent\textbf{결과~\#200} (\textit{마스 첨형식, $R=9.5337$, 홀수}):
$\mathrm{Re}(s)=\tfrac{1}{2}$, $t\in[3,50]$에서 19개의 영점.

\begin{itemize}
  \item P1 (FE): $|\Lambda(s)-\Lambda(1-s)|/|\Lambda(s)|=0.00$ (5개 검사점,
        \texttt{dps}$=80$).
  \item P2 (영점 수): 부호변화 이분법으로 19개 확인
        ($\approx 43$분; $\rho_1=0.5+3.264i$, \ldots, $\rho_{19}=0.5+48.865i$).
  \item P3 ($\kappa\delta^{2}$): 5영점 로그-로그 기울기
        $\mathbf{2.0003\pm 0.0003}$, $R^{2}=1.000000$ (전영점).
        표~\ref{tab:maass_kd2}에 각 영점의 상세를 수록한다.
  \item P4 (모노드로미): $5/5=2.0000\pi$ (등고선 적분, 반지름 $0.5$,
        64단계).
\end{itemize}
종합: \textbf{5/5 통과}.  P5 ($\sigma$-유일성): 7개 시험값
$\sigma\in\{0.1,0.2,0.3,0.5,0.7,0.8,0.9\}$ 모두 부호변화 12회로 동일하여
$\sigma=0.5$ 최대 확인.  무게 12, 20의 정칙 고유형식과 달리 B-01 초과가
관측되지 않았는데, 이는 마스 형식의 넓은 영점 간격($\Delta t\approx 3.3$)에
기인한 것으로 보인다.

\begin{table}[ht]
\centering
\caption{마스 최저 5영점의 $\kappa\delta^{2}$ 로그-로그 기울기
  ($R=9.5337$, $\delta\in[0.005,0.1]$, 8값). 전체 $R^{2}=1.000000$.}
\label{tab:maass_kd2}
\small
\begin{tabular}{llrl}
\toprule
영점 & $t$ & 기울기 & 등급 \\
\midrule
$\rho_1$ & $3.264$ & $1.9999$ & $\bigstar\bigstar\bigstar$ \\
$\rho_2$ & $12.418$ & $2.0002$ & $\bigstar\bigstar\bigstar$ \\
$\rho_3$ & $17.577$ & $2.0004$ & $\bigstar\bigstar\bigstar$ \\
$\rho_4$ & $20.674$ & $2.0005$ & $\bigstar\bigstar\bigstar$ \\
$\rho_5$ & $23.553$ & $2.0007$ & $\bigstar\bigstar\bigstar$ \\
\midrule
평균 & & $2.0003\pm 0.0003$ & \\
\bottomrule
\end{tabular}
\end{table}

기울기 $2.0003\pm 0.0003$은 제1--3논문에 걸쳐 검사된 모든 $L$-함수 중
가장 정밀한 P3 측정이다. 예외적 정밀도는 넓은 영점 간격에서 비롯된다:
최소 간격은 $\rho_4$--$\rho_5$의 $\Delta t=2.88$으로,
비슷한 높이에서의 전형적인 정칙 영점 간격보다 훨씬 크다.
따라서 무게-$16$, 무게-$20$ 형식에서 주된 오차 원인인
근접 영점 간섭이 무시할 수 있다.

\subsection{짝수 마스 첨형식 ($R=13.7798$, 결과~\#201)}
\label{ssec:maass_even}

결과~\#200은 첫 번째 \emph{홀수} 마스 첨형식을 검사하였다.
대칭류와 스펙트럼 매개변수로부터의 독립성을 검증하기 위해,
$\mathrm{SL}_2(\mathbb{Z})\backslash\mathbb{H}$ 위의
첫 번째 \emph{짝수} 마스 첨형식을 검사한다.
스펙트럼 매개변수 $R=13.7797513518907389\ldots$,
$\varepsilon=+1$ (짝수 대칭류).
완성 $L$-함수:
\[
  \Lambda(s,u_R^{+}) = \pi^{-s}\,
  \Gamma\!\Bigl(\frac{s+iR}{2}\Bigr)
  \Gamma\!\Bigl(\frac{s-iR}{2}\Bigr)\,L(s,u_R^{+}), \qquad
  \Lambda(s)=\Lambda(1-s), \quad \varepsilon=+1.
\]

\medskip
\noindent\textbf{결과~\#201} (\textit{짝수 마스 첨형식, $R=13.7798$}):
\begin{itemize}
  \item P1 (FE): $|\Lambda(s)-\Lambda(1-s)|/|\Lambda(s)|=0.00$ (5개 검사점,
        \texttt{dps}$=80$).
  \item P3 ($\kappa\delta^{2}$): 5영점 로그-로그 기울기
        $\mathbf{1.9999\pm 0.0008}$, $R^{2}=1.000000$ (전영점).
  \item P4 (모노드로미): $5/5=2.0000\pi$.
  \item P5 ($\sigma$-유일성): 실패 (B-01);
        $\sigma=0.5\to 11$개, $\sigma=0.1,0.2\to 12$개.
        모든 $\mathrm{GL}(2)$ 결과와 일치.
\end{itemize}
종합: \textbf{4/5 통과} ($\bigstar\bigstar\bigstar$).
결과~\#200과 합산하면, 기울기~$2$는 짝수 및 홀수 대칭류 양쪽과
두 개의 서로 다른 스펙트럼 매개변수 ($R=9.53$, $13.78$)에 걸쳐
지속되어, $\kappa\delta^{2}$ 보편성이 마스 스펙트럼 데이터에
무관함을 확인한다.


\subsection{차수 확장: $\mathrm{GL}(3)$과 $\mathrm{GL}(4)$}
\label{ssec:degree_ext}

앞의 모든 $\kappa\delta^{2}$ 측정은 차수-$2$ $L$-함수
(GL(2) 위의 자기형식)에 관한 것이다.
핵심 질문은 기울기~$2$가 셀베르그 부류의 고차수 $L$-함수에서도
지속되는가이다. 대칭 거듭제곱 $L$-함수를 통해 차수~$3$과 차수~$4$를 검사한다.

\paragraph{$\mathrm{GL}(3)$: $\mathrm{sym}^{2}(E, \text{11a1})$
(결과~\#202).}
\label{par:gl3sym2}
타원곡선 11a1의 대칭제곱 $L$-함수는
차수-$3$, 도체 $N=121$인 자기형식 $L$-함수로서
$\Lambda(s)=\Lambda(3-s)$, $\sigma_{\mathrm{crit}}=3/2$이다.
PARI의 \texttt{lfunzeros}로 발견된 약 30개의 영점 중
다섯 개의 영점 ($t\in[3.9,12.1]$, 간격 $>1.0$)을 선택한다.

\medskip
\noindent\textbf{결과~\#202}: $\kappa\delta^{2}$ 로그-로그 기울기
$\mathbf{2.0000\pm 0.0001}$, $R^{2}=1.000000$ (전영점).
모노드로미 $5/5=2.0000\pi$.
$\sigma$-유일성: $5/5$ 통과
(중심 $\sigma=1.5$가 최대---$\mathrm{GL}(2)$와 달리 차수~$3$에서는
B-01 초과 없음).
FE: $|\Lambda(s)-\Lambda(3-s)|/|\Lambda(s)|=0.00$.
종합: \textbf{5/5 통과} ($\bigstar\bigstar\bigstar$).

$\pm 0.0001$의 오차바는 최우수 $\mathrm{GL}(2)$ 측정
(\#200, 마스 홀수)과 동일하다. 이는 문헌 최초의 차수-$3$
$\kappa\delta^{2}$ 기울기 측정이다.

\paragraph{$\mathrm{GL}(4)$: $\mathrm{sym}^{3}(\Delta)$
(결과~\#203).}
\label{par:gl4sym3}
라마누잔 $\Delta$-함수의 대칭세제곱 $L$-함수는
차수-$4$, 도체 $N=1$, 동기적 무게~$33$인 자기형식 $L$-함수로서
$\Lambda(s)=\Lambda(1-s)$, $\sigma_{\mathrm{crit}}=1/2$이다
(PARI 정규화, $\gamma_{V}=[5.5,6.5,16.5,17.5]$).
$L(1/2)\approx 0$ (근 수 $\varepsilon=-1$).
다섯 개의 영점 ($t\in[4.2,12.1]$, 간격 $>1.0$)은
$|L(1/2+it)|$ 최소점 수동 탐색으로 발견한다
(황금분할 정밀화; 큰 감마 이동에 대해 \texttt{lfunzeros}가
신뢰할 수 없으므로).

\medskip
\noindent\textbf{결과~\#203}: $\kappa\delta^{2}$ 로그-로그 기울기
$\mathbf{2.0000\pm 0.0001}$, $R^{2}=1.000000$ (전영점).
모노드로미 $5/5=2.0000\pi$.
$\sigma$-유일성: 실패 (중심 계수~$=24$, $\sigma=0.40$에서 최대 계수~$=30$),
차수-${\geq}2$ 구조적 실패 패턴 (B-05)과 일치.
FE: \texttt{lfuncheckfeq}~$=-51$ 자릿수.
종합: \textbf{4/5 통과} ($\bigstar\bigstar\bigstar$).

개별 영점 기울기는 $1.9999$, $2.0000$, $2.0001$, $1.9999$,
$2.0000$---각각 이론값~$2$와 구별 불가능하다.
GL(3) 결과와 합산하면, 기울기~$2$는 이제 셀베르그 부류
차수 $\{1,2,3,4\}$에 걸쳐 확립된다.

\paragraph{$\mathrm{GL}(5)$: $\mathrm{sym}^{4}(\text{11a1})$
(결과~\#205).}
\label{par:gl5sym4}
차수 커버리지를 $5$까지 확장하기 위해 대칭네제곱
$L(s,\mathrm{sym}^4(\text{11a1}))$을 검사한다: 도체 $N=14641=11^4$,
차수~$5$, 산술 정규화 $\sigma_{\mathrm{crit}}=k/2=2.5$,
$\gamma_V=[-2,-1,0,0,1]$ (PARI \texttt{lfunsympow(11a1,4)}).
FE: \texttt{lfuncheckfeq}~$=-116$ 자릿수.

\medskip
\noindent\textbf{결과~\#205}: $\kappa\delta^{2}$ 로그-로그 기울기
$\mathbf{1.9999\pm 0.0004}$, $R^{2}=1.000000$ (전영점, $t\in[1.49, 7.35]$,
간격 $>1.0$).
모노드로미 $5/5=2.0000\pi$.
$\sigma$-유일성: 실패 (중심=36, 최대=68),
차수-${\geq}2$ 구조적 패턴 (B-05)과 일치.
종합: \textbf{4/5 통과} ($\bigstar\bigstar\bigstar$).

개별 기울기: $2.0000$, $1.9996$, $1.9995$, $2.0000$, $2.0006$.
이로써 기울기~$2$ 보편성이 차���~$5$에서 확인되어,
셀베르그 부류 차수 $\{1,2,3,4,5\}$를 완전히 커버한다.
제1논문의 결과~\#85에서 이미 GL(5)에서의 고전적 4성질을 검증하였으며,
이번 측정은 다중-$\delta$ 정밀도의 정량적 $\kappa\delta^{2}$ 기울기를 추가한다.
차수~$5$는 계산적 제약 ($\mathrm{sym}^4(\text{37a1})$은 $N=37^4=1{,}874{,}161$)으로
하나의 $L$-함수에 한정된다.

\textcolor{ForestGreen}{[결과 \#205, 확인 2026-04-21;
sym$^4$(11a1): $N{=}14641$, $k{=}5$, $\gamma_V{=}[-2,-1,0,0,1]$;
FE${=}-116$ 자릿수; 기울기${=}1.9999{\pm}0.0004$; 모노 $5/5{=}2\pi$;
$\sigma$-유일성 실패 (구조적); 12행 비교표; 차수 1--5 완성]}


\subsection{$\kappa\delta^{2}$의 무게, 스펙트럼 및 차수 불변성}
\label{ssec:weightinv}

표~\ref{tab:weightuniv}은 지금까지 검사된 모든 모듈러 무게와 스펙트럼 매개변수에
걸친 $\kappa\delta^{2}$ 로그-로그 기울기
($\log(\kappa\delta^{2}-1)$ 대 $\log\delta$의 기울기)를 수집한다.

\begin{table}[ht]
\centering
\caption{$\kappa\delta^{2}$ 로그-로그 기울기의 무게, 스펙트럼, 차수 및 가족 불변성.
  차수 $1$--$6$, 무게 $k\in\{0,1,2,12,16,20\}$, 마스 스펙트럼 매개변수
  $R=9.53,13.78$, 대칭 거듭제곱 $\mathrm{sym}^{2}$--$\mathrm{sym}^{5}$,
  두 개의 아틴 $L$-함수 ($S_3$, $S_5$), 하나의 Rankin--Selberg
  합성곱 ($\text{11a1}\times\text{37a1}$),
  하나의 자기쌍대 디리클레 $L$-함수 ($\chi_{-7}$), 세 개의 \emph{비자기쌍대}
  복소 지표 ($\chi\bmod 5,7,13$; 결과~\#216)를 아우르는
  열아홉 개의 데이터 점. 모든 $\sigma$-방향 측정은 다중-$\delta$ 로그-로그 프로토콜을 사용한다.}
\label{tab:weightuniv}
\small
\begin{tabular}{llllrrl}
\toprule
유형 & $L$-함수 & 차수 & $\sigma_{\mathrm{crit}}$ & 기울기 & $\pm\sigma$ & 등급 \\
\midrule
$k{=}0$  & 리만 $\zeta$ & $\mathrm{GL}(1)$ & $\tfrac{1}{2}$ & $\approx 2.0$ & --- & $\bigstar\bigstar\bigstar$ \\
디리클레 & $L(s,\chi_{-7})$ (\#213) & $\mathrm{GL}(1)$ & $\tfrac{1}{2}$ & $\mathbf{2.0000}$ & $\mathbf{0.0014}$ & $\bigstar\bigstar\bigstar$ \\
$k{=}1$  & 아틴 $S_3$ (\#207) & $\mathrm{GL}(2)$ & $\tfrac{1}{2}$ & $\mathbf{2.0000}$ & $\mathbf{0.0000}$ & $\bigstar\bigstar\bigstar$ \\
$k{=}2$  & EC 11a1 & $\mathrm{GL}(2)$ & $1$ & $\approx 2.0$ & --- & $\bigstar\bigstar\bigstar$ \\
마스 홀수 & $u_{R}^{-}$ ($R{=}9.53$) & $\mathrm{GL}(2)$ & $\tfrac{1}{2}$ & $\mathbf{2.0003}$ & $\mathbf{0.0003}$ & $\bigstar\bigstar\bigstar$ \\
마스 짝수 & $u_{R}^{+}$ ($R{=}13.78$) & $\mathrm{GL}(2)$ & $\tfrac{1}{2}$ & $1.9999$ & $0.0008$ & $\bigstar\bigstar\bigstar$ \\
$k{=}12$ & 라마누잔 $\Delta$ & $\mathrm{GL}(2)$ & $6$ & $2.0008$ & $0.0006$ & $\bigstar\bigstar\bigstar$ \\
$k{=}16$ & $\Delta\cdot E_4$ & $\mathrm{GL}(2)$ & $8$ & $1.9989$ & $0.0035$ & $\bigstar\bigstar\bigstar$ \\
$k{=}20$ & $\Delta\cdot E_8$ & $\mathrm{GL}(2)$ & $10$ & $1.9984$ & $0.0055$ & $\bigstar\bigstar\bigstar$ \\
$\mathrm{sym}^{2}$ & $\mathrm{sym}^{2}(\text{11a1})$ & $\mathrm{GL}(3)$ & $\tfrac{3}{2}$ & $\mathbf{2.0000}$ & $\mathbf{0.0001}$ & $\bigstar\bigstar\bigstar$ \\
$\mathrm{sym}^{3}$ & $\mathrm{sym}^{3}(\Delta)$ & $\mathrm{GL}(4)$ & $\tfrac{1}{2}$ & $\mathbf{2.0000}$ & $\mathbf{0.0001}$ & $\bigstar\bigstar\bigstar$ \\
$\mathrm{sym}^{3}$ & $\mathrm{sym}^{3}(\text{37a1})$ & $\mathrm{GL}(4)$ & $2$ & $\mathbf{1.9999}$ & $\mathbf{0.0005}$ & $\bigstar\bigstar\bigstar$ \\
아틴 $S_5$ & $\rho_4$ (\#209) & $\mathrm{GL}(4)$ & $\tfrac{1}{2}$ & $\mathbf{1.9999}$ & $\mathbf{0.0003}$ & $\bigstar\bigstar\bigstar$ \\
R-S & $\text{11a1}\times\text{37a1}$ (\#210) & $\mathrm{GL}(4)$ & $1$ & $\mathbf{1.9960}$ & $\mathbf{0.0072}$ & $\bigstar\bigstar\bigstar$ \\
$\mathrm{sym}^{4}$ & $\mathrm{sym}^{4}(\text{11a1})$ & $\mathrm{GL}(5)$ & $\tfrac{5}{2}$ & $\mathbf{1.9999}$ & $\mathbf{0.0004}$ & $\bigstar\bigstar\bigstar$ \\
$\mathrm{sym}^{5}$ & $\mathrm{sym}^{5}(\text{11a1})$ & $\mathrm{GL}(6)$ & $3$ & $\mathbf{1.9980}$ & $\mathbf{0.0036}$ & $\bigstar\bigstar\bigstar$ \\
\midrule
\multicolumn{7}{l}{\textit{비자기쌍대 복소 지표 (결과~\#216)}} \\
$\chi_{5}$ & $\chi\bmod 5$ 차수~4 & $\mathrm{GL}(1)$ & $\tfrac{1}{2}$ & $\mathbf{2.0005}$ & $\mathbf{0.0022}$ & $\bigstar\bigstar\bigstar\bigstar$ \\
$\chi_{7}$ & $\chi\bmod 7$ 차수~6 & $\mathrm{GL}(1)$ & $\tfrac{1}{2}$ & $\mathbf{1.9997}$ & $\mathbf{0.0020}$ & $\bigstar\bigstar\bigstar\bigstar$ \\
$\chi_{13}$ & $\chi\bmod 13$ 차수~12 & $\mathrm{GL}(1)$ & $\tfrac{1}{2}$ & $\mathbf{2.0009}$ & $\mathbf{0.0035}$ & $\bigstar\bigstar\bigstar\bigstar$ \\
\bottomrule
\multicolumn{7}{l}{\small 관습 A (전체): log$(\kappa\delta^2-1)$ 대 log$(\delta)$, 이론적 기울기 $2$.} \\
\end{tabular}
\end{table}

정칙 형식의 경우 $k=12$, $16$, $20$의 기울기는 일치한다: 3-무게 평균은
$1.9994\pm 0.0010$이며, 최대 쌍별 편차는 $0.0024$이다.
두 마스 형식 ($R=9.5337$ 홀수, $R=13.7798$ 짝수)은
질적으로 다른 데이터 점---비정칙, 무게~$0$, 스펙트럼 감마 인자---을 추가하되,
그 기울기 ($2.0003\pm 0.0003$, $1.9999\pm 0.0008$)는
가장 정밀한 측정들에 속한다.

대칭 거듭제곱을 통한 차수 확장은 동등하게 인상적이다:
차수~$3$의 $\mathrm{sym}^{2}(\text{11a1})$과
차수~$4$의 $\mathrm{sym}^{3}(\Delta)$ 모두 기울기 $2.0000\pm 0.0001$을 산출하여,
이론값에 소수점 네 자리까지 일치한다.
아틴 $S_5$ 표준표현 $\rho_4$ (결과~\#209: 비가환 차수-4, 도체~$2869$,
기울기~$1.9999\pm 0.0003$)는 대칭 거듭제곱 계열을 넘어선 보편성을 확인한다.
Rankin--Selberg 합성곱 $L(s,\text{11a1}\times\text{37a1})$
(결과~\#210: 차수~4, $N=407$, 기울기~$1.9960\pm 0.0072$)은
최초의 텐서 곱 검증을 제공하여 네 번째 독립적 차수-4 구성을 확립한다.
제1논문의 $k=0$ (리만 $\zeta$)과 $k=2$ (EC 11a1) 결과와 합산하면,
기울기는 셀베르그 부류 차수 $\{1,2,3,4,5,6\}$,
정칙 무게 $k\in\{0,1,2,12,16,20\}$, 두 개의 비정칙 마스 형식,
네 개의 독립적 GL(4) $L$-함수 (대칭 거듭제곱 2개, 아틴, 랭킨-셀베르그),
GL(5) $L$-함수 하나, GL(6) $L$-함수 하나를 아우르는
열아홉 개의 데이터 점에 걸쳐 $2.0\pm 0.008$이다.

\begin{theorem}[기울기 보편성]
\label{thm:slopeuniv}
  완비 $L$-함수 $\Lambda(s)$가 다음을 만족한다고 하자:
  \begin{enumerate}
    \item[\textup{(FE)}] $\Lambda(s) = \varepsilon\,\widetilde{\Lambda}(k-s)$
          ($k\in\mathbb{R}$, $|\varepsilon|=1$, $\widetilde{\Lambda}$는
          대조쌍\textup{(}자기쌍대이면 $\widetilde{\Lambda}=\Lambda$\textup{)});
          그리고
    \item[\textup{(SR)}] $\widetilde{\Lambda}(\bar{s})
          = \overline{\Lambda(s)}$
          \textup{(}슈바르츠 반사\textup{)}.
  \end{enumerate}
  두 조건은 모든 보형 $L$-함수에서 성립하며, 특히 자기쌍대가 \emph{불필요}하다.
  임계선 위의 단순 영점 $\rho=\tfrac{k}{2}+i\gamma$에서
  로랑 전개를
  $\Lambda'/\Lambda(s) = (s-\rho)^{-1} + c_0 + c_1(s-\rho) + \cdots$라 쓰면:
  \begin{enumerate}
    \item[\textup{(i)}] $\operatorname{Re}(c_0)=0$ \textup{(}$c_0$는 순허수\textup{)};
    \item[\textup{(ii)}] $\operatorname{Im}(c_1)=0$ \textup{(}$c_1$은 실수\textup{)};
    \item[\textup{(iii)}] $\kappa(\rho{+}\delta)\,\delta^{2}
          = 1 + A\,\delta^{2} + O(\delta^{3})$,
          여기서 $A = \operatorname{Im}(c_0)^{2} + 2c_1$.
  \end{enumerate}
  특히, $\log(\kappa\delta^{2}-1) = \log A + 2\log\delta + O(\delta)$이므로
  로그-로그 기울기는 정확히~$2$이다.
\end{theorem}

\begin{proof}
  $s=\rho+\delta$ ($\delta\in\mathbb{R}$)에서:
  $\Lambda'/\Lambda = 1/\delta + c_0 + c_1\delta + \cdots$.
  (FE)에서 $\Lambda'/\Lambda(s)=-\widetilde{\Lambda}'/\widetilde{\Lambda}(k{-}s)$.
  $k-s=\bar{\rho}-\delta$이고 $\bar{\rho}$는 $\widetilde{\Lambda}$의 영점이므로
  $\widetilde{\Lambda}'/\widetilde{\Lambda}(\bar{\rho}{-}\delta)
  = -1/\delta + \tilde{c}_0 - \tilde{c}_1\delta + \cdots$.
  (SR)에서 $s'=\rho-\delta$ (따라서 $\bar{s}'=\bar{\rho}-\delta$)로 놓으면:
  $\widetilde{\Lambda}'/\widetilde{\Lambda}(\bar{\rho}{-}\delta)
  = \overline{\Lambda'/\Lambda(\rho{-}\delta)}
  = -1/\delta + \bar{c}_0 - \bar{c}_1\delta + \cdots$.
  (FE)에 대입하면:
  \[
    \tfrac{1}{\delta}+c_0+c_1\delta
    = -(-\tfrac{1}{\delta}+\bar{c}_0-\bar{c}_1\delta)
    = \tfrac{1}{\delta}-\bar{c}_0+\bar{c}_1\delta.
  \]
  계수 비교: $c_0=-\bar{c}_0$ (따라서 $\operatorname{Re}(c_0)=0$)이고
  $c_1=\bar{c}_1$ (따라서 $\operatorname{Im}(c_1)=0$).
  그러면 $\kappa\delta^{2}=|1+c_0\delta+c_1\delta^2|^2
  =1+2\operatorname{Re}(c_0)\delta+(|c_0|^2+2\operatorname{Re}(c_1))\delta^2+O(\delta^3)
  =1+A\delta^2+O(\delta^3)$
  ($\operatorname{Re}(c_0)=0$이 $\delta^1$ 항을 소멸시키고,
  $|c_0|^2=\operatorname{Im}(c_0)^2$, $\operatorname{Re}(c_1)=c_1$).
\end{proof}

\begin{corollary}[진폭 공식]
\label{cor:amplitude}
  차선도 계수 $A=\operatorname{Im}(c_0)^2+2c_1$은
  \emph{국소 밀도} 기여 $\operatorname{Im}(c_0)^2$
  (인접 영점을 통해 $c_0=\sum_{\rho'\neq\rho}(\rho-\rho')^{-1}$)와
  \emph{감마 인자} 기여 $2c_1$ (완비 $L$-함수의 디감마 항)으로 분해된다.
\end{corollary}

\noindent\textit{수치 검증} (결과~\#214, B-34):
자기쌍대 $L$-함수 4개($\zeta$, $L(s,\chi_{-3})$, $L(s,\chi_{-7})$,
$L(s,\chi_{-11})$)의 $20$개 영점에서
$|\operatorname{Re}(c_0)|=0$ (150자리 산술에서 기계 정밀도),
평균 기울기 $2.0003\pm 0.0022$, 평균 $A$-예측 오차 $0.5\%$ (최대 $1.6\%$).
결과~\#216는 \emph{비자기쌍대} 복소 지표 3개 ($\chi\bmod 5$ 차수~4,
$\chi\bmod 7$ 차수~6, $\chi\bmod 13$ 차수~12)의 $15$개 영점으로 확장하여,
모두 $|\operatorname{Re}(c_0)|=0$, 평균 기울기 $2.0004\pm 0.0024$를 확인했다---정리가 비자기쌍대 경우까지 포괄함을 증명한다.
이 정리는 기존 정리~\ref{thm:rankinv} (명시적 $A$ 공식 없이
$\kappa\delta^2=1+O(\delta^2)$를 확립)를 포괄하며,
표~\ref{tab:weightuniv}의 기울기~$2$ 관찰을 수치적 규칙성에서
수학적 정리로 승격시킨다.
셀베르그 부류 차수 $\{1,2,3,4,5,6\}$에 걸친 19개 데이터 포인트가
예측된 기울기~$2$를 $\pm 0.008$ 내에서 확인하여
\emph{구성 보편성}을 확립한다.

\paragraph{진폭 스케일링: 차선도 계수 $A(t_{0})$.}
기울기는 보편 상수($=2$)이지만, 진폭
$A(t_{0}) \;\coloneqq\; \kappa(\sigma_{\mathrm{crit}}+\delta,\,t_{0}) - 1/\delta^{2}$은
$\delta\to 0$에서 영점과 $L$-함수에 따라 달라진다.
$\delta=0.01$에서 차수 $1$--$6$에 걸친 $52$개 영점으로부터 $A$를 추출하면
(결과~\#212), 2변수 의존성이 관찰된다:
\begin{equation}\label{eq:Ascaling}
  A(d,N) \;\approx\; \alpha\,d^{2} \;+\; \beta\,\log N \;+\; \gamma,
  \qquad \alpha\approx 0.39,\;\beta\approx 2.19,\;\gamma\approx 0.10,
\end{equation}
표본 내 $R^{2}=0.983$이다. $d^{2}$ 항은 아다마르 곱 분해
$A = B(t_{0})^{2} + 2H_{1}(t_{0})$에서 유래하며,
여기서 $H_{1}=\sum_{\rho'\neq\rho}(\mathrm{Im}\,\rho'-t_{0})^{-2}$은
영점 밀도(따라서 $d$)와 함께 증가하고, $\log N$ 항은 완비
$L$-함수의 감마 인자 스털링 근사에서 발생한다.

\emph{표본 외 검증.}  아틴 $S_5$ $L$-함수
(결과~\#209: $d=4$, $N=2869$)에 대해 식~\eqref{eq:Ascaling}은
$A\approx 23.9$를 예측하며, 5개 영점의 측정 평균은 $24.9$이다 (오차~$4\%$).
$L(s,\text{11a1}\times\text{37a1})$ (결과~\#210: $d=4$, $N=407$)에서는
예측값 $19.6$ 대 측정값~$30.7$이다 (인접 영점 쌍으로 인한 높은 분산).
정성적 경향---$A$가 $d^{2}$과 $\log N$ 모두에 따라 증가---은 강건하며,
정확한 계수는 안정화를 위해 추가 데이터가 필요하다.
\textcolor{ForestGreen}{[결과 \#212, 확인 2026-04-21;
6차수, 52영점; 아틴 $S_5$ 표본 외 검증 오차 $4\%$;
관측: $A(d,N)\propto d^{2}+\log N$ (잠정 계수)]}

\emph{$A(t_0)$의 내재적 변동성.}  $\zeta(s)$의 $269$개 영점에 대해
아다마르 합으로 $H_1$을 직접 계산하면 (결과~\#215), 변동계수
$\mathrm{CV}(H_1)\approx 42\%$가 관찰된다. 이는 국소 영점 간격의
GUE형 요동에서 비롯된다.  정규화된 비 $H_1\cdot\Delta(t_0)^2$
($\Delta=2\pi/\log(t_0/2\pi)$는 평균 간격)의 평균은 $t_0\in[50,500]$에서
$4.58\pm 1.95$로 안정적이나, 점별 변동이 커서 정밀한 보편 상수를 추출하기
어렵다.  $B(t_0)^2$ 항 (힐베르트형 합
$\sum_{\rho'\neq\rho}(\mathrm{Im}\,\rho'-t_0)^{-1}$에서 유래)은 $A$의
${\approx}24\%$를 기여하며, 평균 0, 표준편차 ${\approx}1.0$으로 요동한다.
이 내재적 $O(1)$ 변동성은 회귀식~\eqref{eq:Ascaling}이 표본 내
$R^2=0.98$이면서도 점 예측에서 ``잠정적''인 이유를 설명한다.
\textcolor{ForestGreen}{[결과 \#215, 확인 2026-04-21;
$\zeta$의 269영점; $H_1=\sum 1/(\gamma_0-\gamma_n)^2$ PARI 대비
$0.7\%$ 검증; CV$=42\%$; $2H_1$이 $A$의 $76\%$ 기여]}


% ============================================================================
\section{논의}
\label{sec:disc}
% ============================================================================

\paragraph{보편성.}
네 가지 $\xi$-다발 성질은 확장된 $L$-함수 족에 걸쳐 통과한다:
타원곡선 $L$-함수 (계수 0--3, 도체 11--10099, 8개 곡선, $\varepsilon=\pm 1$),
랭킨-셀베르그 $\mathrm{GL}(4)$ (2쌍),
데데킨트 $\zeta$-함수 (7개의 허수 이차체),
유수 2--3인 3개의 비-오일러 형태를 포함한 엡스타인 $\zeta$-함수.
제1논문의 $\mathrm{GL}(1)$--$\mathrm{GL}(5)$ 결과 (81개의 수치 실험)와 합산하면,
프레임워크는 이제 $L$-함수의 주요 대수 족 전체를 포괄한다.

\paragraph{무게, 스펙트럼, 차수 및 가족 보편성.}
결과~\#125--\#127, \#200--\#201, \#202--\#203은 $\kappa\delta^{2}$ 로그-로그
기울기가 모듈러 무게, 정칙성, 마스 대칭류, 셀베르그 부류 차수에 걸쳐
불변임을 확립한다.
정칙 고유형식의 3-무게 평균은 $1.9994\pm 0.0010$ ($k=12,16,20$)이며;
두 마스 첨형식은 $2.0003\pm 0.0003$ ($R=9.53$, 홀수)과
$1.9999\pm 0.0008$ ($R=13.78$, 짝수)을 산출한다.
대칭 거듭제곱을 통한 차수 확장은 가장 정밀한 측정을 제공한다:
$\mathrm{GL}(3)$의 $\mathrm{sym}^{2}(\text{11a1})$과
$\mathrm{GL}(4)$의 $\mathrm{sym}^{3}(\Delta)$ 모두 기울기
$2.0000\pm 0.0001$을 산출한다.
제1논문의 $k=0,2$ 결과, 아틴 결과 ($k=1$), Rankin--Selberg 텐서 곱
(결과~\#210, 기울기 $1.9960\pm 0.0072$), 네 개의 원시 디리클레 지표
(결과~\#213, 기울기 $2.0002\pm 0.0001$)를 합산하면,
기울기는 셀베르그 부류 차수 $\{1,2,3,4,5,6\}$와
일곱 개의 구성 가족에 걸쳐 $2.0\pm 0.008$이다:
열아홉 개의 데이터 점.
정리~\ref{thm:slopeuniv}는 이를 해석적으로 증명한다:
함수방정식과 슈바르츠 반사가 로랑 전개에서 $\operatorname{Re}(c_0)=0$을
강제하여 $\delta^1$ 항을 소멸시키므로,
기울기~$2$는 비자기쌍대를 포함한 \emph{모든} $L$-함수의 수학적 필연이다.
따름정리~\ref{cor:amplitude}에 의해, $c_{1}=\sum(\gamma-\gamma')^{-2}$ 항이
$A$의 약 $76\%$를 지배하며, 나머지는 국소 영점 밀도
$\mathrm{Im}(c_0)^2$에서 기인한다.

\paragraph{오일러 곱의 불필요성.}
정리~\ref{thm:feonly}와 결과~\#114는 오일러 곱이
$\xi$-다발 성질에 필요하지 않음을 확립한다.
이는 $\xi$-다발 기하가 함수방정식 구조만의 귀결이라는
해석과 부합한다.

\paragraph{판별 정리.}
정리~\ref{thm:disc}는 $\kappa\delta^{2}$에 대한
$O(\delta)$ 보정항을 임계선으로부터의 영점의 위치와
연결하는 최초의 \emph{해석적} 진술을 제공한다.
수치 검증 (결과~\#115--\#117)은 유일하게 알려진 명시적 시험 사례
(DH 함수)에 대해 이 연결을 확인한다.
리만 $\zeta$-함수에 대해서는 임계선 외 영점이 알려져 있지 않으며;
정리~\ref{thm:disc}는 그러한 영점이 존재한다면
감지 가능한 $O(\delta)$ 신호를 생성할 것임을 함의한다.

\paragraph{한계.}
점근식 $c_{1}\cdot d \to 1$은 $N=4$개의 데이터 점으로 검증되었다.
체계적인 $\delta$-스윕 (결과~\#116)은
$N=3$개의 임계선 외 영점에 한정된다
(OFF\#4는 $\delta<0.1$에서 정밀도 불충분).
$t\in[200,400]$ 범위에서는 새로운 임계선 외 영점이
발견되지 않았다 (22개의 후보점 검사, 0개 확인),
DH 임계선 외 영점의 희소성과 일치한다.

\paragraph{향후 방향.}
\begin{enumerate}
  \item \textit{힐베르트 모듈러 형식}: 스펙트럼 보편성을 차수 $>1$인 수체로 확장.
  \item \textit{고차수 확장}: 차수~$7+$를 $\mathrm{sym}^{6}$ 또는 기타
        자기형식 구성을 통해 달성.
  \item \textit{셀베르그 부류 일반화}: 정리~\ref{thm:disc}는
        FE만 사용한다; 셀베르그 부류 차수에 의한 $c_{1}$을
        제시하도록 정밀화할 수 있는가?
  \item \textit{$\sigma$-국소화}: 비수평 오프셋
        (복소 $\delta$)으로의 판별 확장.
  \item \textit{추가 DH 임계선 외 영점}: $t>400$에 대한
        더 높은 정밀도에서의 탐색 확장.
\end{enumerate}


% ============================================================================
\section{결론}
\label{sec:conclusion}
% ============================================================================

Kim~(2026)의 $\xi$-다발 프레임워크를 네 가지 방향으로 확장하였다.
첫째, 보편성: 열한 개의 새로운 수치 결과
(\#107--\#117)가 타원곡선, 랭킨-셀베르그, 데데킨트, 엡스타인
$L$-함수에 걸쳐 네 가지 $\xi$-다발 성질을 확인하며;
특히 비-오일러 엡스타인 $\zeta$-함수 (정리~6)는
오일러 곱이 필요하지 않음을 보여준다.
둘째, 판별: 정리~7은 $\kappa\delta^{2}$의 첫 번째 테일러 계수
$c_{1}$이 영점이 임계선 위에 있을 때에만 소멸함을 해석적으로 증명하며;
$\delta$-스윕과 직접 해석 계산을 통한 수치 검증이
알려진 네 개의 DH 임계선 외 영점 전체에서 이를 확인한다.
셋째, 무게 및 스펙트럼 보편성: 결과~\#125--\#127과 \#200--\#201이
$\kappa\delta^{2}$ 로그-로그 기울기가 모듈러 무게 ($k=12,16,20$)와
정칙성 (두 마스 형식, $R=9.53$과 $13.78$), 대칭류 (홀수 및 짝수)에
독립적으로 $2.0\pm 0.001$임을 확립하며, 정리~8을 확인한다.
넷째, 차수 및 가족 보편성: 결과~\#202--\#206은 기울기 검증을
$\mathrm{GL}(2)$에서 $\mathrm{GL}(6)$으로 대칭 거듭제곱을 통해 확장하며;
결과~\#209는 비가환 차수-4 아틴 $L$-함수
(표준표현 $S_5$, 기울기 $1.9999\pm 0.0003$)를 추가하고;
결과~\#210은 Rankin--Selberg 텐서 곱
$L(s,\text{11a1}\times\text{37a1})$ (기울기 $1.9960\pm 0.0072$)을 추가하여,
셀베르그 부류 차수 $\{1,2,3,4,5,6\}$에 걸쳐
일곱 개의 구성 가족으로부터 열아홉 개의 독립 데이터 점으로
정리~8을 확인한다;
결과~\#213은 네 개의 원시 디리클레 지표 (도체 $N=3,4,5,7$,
기울기 $2.0002\pm 0.0001$)를 추가하여 여섯 번째 가족을 확립하며;
결과~\#216은 비자기쌍대 복소 지표 3개 ($\chi\bmod 5,7,13$, 기울기
$2.0004\pm 0.0024$)를 추가하여 일곱 번째 가족을 확립한다.

정리~5--8은 하나의 일관된 그림을 형성한다:
임계선 위에서 $\xi$-다발 곡률은 검사된 모든 $L$-함수 족,
모듈러 무게, 스펙트럼 유형에 걸쳐 $\kappa\delta^{2}=1+O(\delta^{2})$를
보편적으로 만족하고; 임계선 외에서는 질적으로 다른 $O(\delta)$ 보정항이
나타나며, 이는 해석적으로 함수방정식에 의한 거울 대칭의 파괴로 추적된다.


% ============================================================================
% 부록
% ============================================================================
\appendix

\section{종합 요약표: 스물여덟 개 결과 전체}
\label{app:summary}

\begin{longtable}{lllp{5.5cm}}
\caption{결과~\#107--\#117, \#125--\#127, \#200--\#216 요약 (동반 논문). 등급:
  \textbf{$\bigstar\bigstar\bigstar$} 강양성,
  \textbf{$\bigstar\bigstar$} 양성, \textbf{C} 조건부.}
\label{tab:summary}\\
\toprule
결과 & 절 & 등급 & 핵심 발견 \\
\midrule
\endfirsthead
\multicolumn{4}{c}{\tablename~\thetable{} (계속)}\\
\toprule
결과 & 절 & 등급 & 핵심 발견 \\
\midrule
\endhead
\bottomrule
\endfoot
\#107 & \S\ref{ssec:epsilon}      & \textbf{$\bigstar\bigstar$}\,C  & EC 8 곡선; $\varepsilon=-1$에서 P4 실패 (Hardy $Z$ 방식); P3 전체 통과 \\
\#108 & \S\ref{ssec:epsilon}      & \textbf{$\bigstar\bigstar\bigstar$}    & EC 8 곡선, $\varepsilon$-보정; 전체 4/4 통과; $\kappa\delta^{2}=1.000$ \\
\#109 & \S\ref{ssec:thm5}         & \textbf{$\bigstar\bigstar$}     & $\kappa_{\text{near}}$ 보편성; $\delta=0.01$에서 CV $<0.13\%$; 계수/도체 무관 \\
\#110 & \S\ref{ssec:mono}         & \textbf{$\bigstar\bigstar\bigstar$}    & EC 8 곡선, 120 영점; $|\Delta\!\arg/\pi|=1.000000$; 등고선 $=2.000000\pi$ \\
\#111 & \S\ref{ssec:conductor}    & \textbf{$\bigstar\bigstar$}     & 도체 스케일링 $\Delta\kappa\approx 2\log N\times 10^{-3}$; $R^{2}=0.73$; 효과 $<2\%$ \\
\#112 & \S\ref{sec:gl4rs}         & \textbf{$\bigstar\bigstar$}     & R-S GL(4), 2쌍; 4/4 통과; $\kappa\delta^{2}=1.014$, $1.033$ \\
\#113 & \S\ref{sec:dedekind}      & \textbf{$\bigstar\bigstar\bigstar$}    & Dedekind $\zeta$, 7 체 (h=1,3); 4/4 통과; 전체 $\kappa\delta^{2}<1.005$ \\
\#114 & \S\ref{sec:epstein}       & \textbf{$\bigstar\bigstar\bigstar$}    & Epstein $\zeta$; 비-오일러 3개 (h=2,3); 4/4 통과; FE-only 충분성 (정리~6) \\
\#115 & \S\ref{ssec:dhtest}       & \textbf{$\bigstar\bigstar\bigstar$}    & DH 임계선 위/외 판별; $1.002\pm0.001$ 대 $1.210\pm0.099$; 완전 분리 \\
\#116 & \S\ref{ssec:sweep}        & \textbf{$\bigstar\bigstar\bigstar$}    & $\delta$-스윕; 기울기 $1.999\pm0.003$ (위) 대 $0.963\pm0.039$ (외); 정리~7 확인 \\
\#117 & \S\ref{ssec:c1}           & \textbf{$\bigstar\bigstar$}     & $c_{1}=\mathrm{Re}(\Lambda''/\Lambda')$ 검증 ($<0.5\%$, $N=4$); $c_{1}\cdot d\to 1$ \\
\#125 & \S\ref{ssec:delta12}      & \textbf{$\bigstar\bigstar\bigstar$}    & 라마누잔 $\Delta$ ($k=12$); 4/5 통과; $\kappa\delta^{2}$ 기울기 $2.0008\pm0.0006$; $\sigma$-유일성 실패 (B-01) \\
\#126 & \S\ref{ssec:w16}          & \textbf{$\bigstar\bigstar\bigstar$}    & 무게-16 첨두형식 ($\Delta\cdot E_4$); 5/5 통과; 기울기 $1.9989\pm0.0035$; 무게 불변성 (정리~8) \\
\#127 & \S\ref{ssec:w20}          & \textbf{$\bigstar\bigstar\bigstar$}    & 무게-20 첨두형식 ($\Delta\cdot E_8$); 4/5 통과; 기울기 $1.9984\pm0.0055$; $\sigma$-유일성 실패 (B-01) \\
\#200 & \S\ref{ssec:maass}        & \textbf{$\bigstar\bigstar\bigstar$}    & 마스 첨형식 ($R{=}9.5337$, 홀수); 5/5 통과; 기울기 $\mathbf{2.0003\pm0.0003}$ \\
\#201 & \S\ref{ssec:maass_even}   & \textbf{$\bigstar\bigstar\bigstar$}    & 짝수 마스 ($R{=}13.7798$); 4/5 통과; 기울기 $1.9999\pm0.0008$; 대칭류 독립성 \\
\#202 & \S\ref{ssec:degree_ext}   & \textbf{$\bigstar\bigstar\bigstar$}    & GL(3) $\mathrm{sym}^{2}(\text{11a1})$; 5/5 통과; 기울기 $\mathbf{2.0000\pm0.0001}$; 최초 차수-3 측정 \\
\#203 & \S\ref{ssec:degree_ext}   & \textbf{$\bigstar\bigstar\bigstar$}    & GL(4) $\mathrm{sym}^{3}(\Delta)$; 4/5 통과; 기울기 $\mathbf{2.0000\pm0.0001}$; 차수 1--4 보편성 \\
\#204 & \S\ref{ssec:degree_ext}   & \textbf{$\bigstar\bigstar\bigstar$}    & GL(4) $\mathrm{sym}^{3}(\text{37a1})$; 4/5 통과; 기울기 $\mathbf{1.9999\pm0.0005}$; 11a1 편향 해소 \\
\#205 & \S\ref{ssec:degree_ext}   & \textbf{$\bigstar\bigstar\bigstar$}    & GL(5) $\mathrm{sym}^{4}(\text{11a1})$; 4/5 통과; 기울기 $\mathbf{1.9999\pm0.0004}$; 차수 1--5 완성 \\
\#206 & \S\ref{ssec:degree_ext}   & \textbf{$\bigstar\bigstar\bigstar$}    & GL(6) $\mathrm{sym}^{5}(\text{11a1})$; 4/5 통과; 기울기 $\mathbf{1.9980\pm0.0036}$; 차수 1--6 완성 \\
\#209 & \S\ref{ssec:degree_ext}   & \textbf{$\bigstar\bigstar\bigstar$}    & 아틴 $S_5$ $\rho_4$; 4/4 통과; 기울기 $\mathbf{1.9999\pm0.0003}$; 독립 차수-4 (제3논문) \\
\#210 & \S\ref{sec:gl4rs}         & \textbf{$\bigstar\bigstar\bigstar$}    & R-S $\text{11a1}\times\text{37a1}$; 실질 4/4 통과; 기울기 $\mathbf{1.9960\pm0.0072}$; 텐서 곱 구성 \\
\#212 & \S\ref{ssec:weightinv}    & \textbf{$\bigstar\bigstar$}           & $A(d,N)$ 스케일링 법칙; $R^{2}=0.983$; 아틴 $S_5$ 표본 외 $4\%$ 오차; 잠정 계수 \\
\#213 & \S\ref{ssec:weightinv}    & \textbf{$\bigstar\bigstar\bigstar$}    & 디리클레 $L(s,\chi_D)$, $D\in\{-3,-4,5,-7\}$; 기울기 $\mathbf{2.0002\pm0.0001}$; 6번째 가족 (도체 독립) \\
\#214 & \S\ref{ssec:weightinv}    & \textbf{$\bigstar\bigstar\bigstar$}    & 기울기 보편성 정리 검증; 4 $L$-함수, 20 영점; $\mathrm{Re}(c_0){=}0$, $\mathrm{Im}(c_1){=}0$; 기울기 $2.0003{\pm}0.0022$; $A$-오차 $0.5\%$ \\
\#215 & \S\ref{ssec:weightinv}    & \textbf{$\bigstar\bigstar$}           & $A(t_0)$ 아다마르 분해; $H_1$ 검증 $0.7\%$; CV$=42\%$ 내재적; $2H_1$이 $A$의 $76\%$ \\
\#216 & \S\ref{ssec:weightinv}    & \textbf{$\bigstar\bigstar\bigstar\bigstar$} & 비자기쌍대 검증: 복소 $\chi$ 3개 (mod 5,7,13), 15 영점, 자기쌍대 대조군 10개; $\mathrm{Re}(c_0){=}0$ (25점 전체); 기울기 $2.0004{\pm}0.0024$; 정리 일반화: 자기쌍대 불필요 \\
\end{longtable}


\section{수치 데이터: DH 임계선 외 영점}
\label{app:dhzeros}

표~\ref{tab:dhdata}는 충분한 정밀도를 갖는
세 개의 DH 임계선 외 영점에 대한 완전한 $\delta$-스윕 데이터를 수록한다
(OFF\#4는 $\delta\geq 0.10$에서만 포함).

\begin{table}[ht]
\centering
\caption{결과~\#116의 $\kappa\delta^{2}$ 값, DH 임계선 외 영점.
  임계선 위 앙상블 평균 (10개 영점)이 비교를 위해 제시된다.}
\label{tab:dhdata}
\small
\begin{tabular}{lrrrrrrrr}
\toprule
영점 & $\delta{=}0.005$ & $0.01$ & $0.015$ & $0.02$ & $0.03$ & $0.05$ & $0.07$ & $0.10$ \\
\midrule
위 (평균) & 1.000048 & 1.000192 & 1.000432 & 1.000767 & 1.001726 & 1.004795 & 1.009399 & 1.019182 \\
OFF\#1 & 1.018828 & 1.037696 & 1.056606 & 1.075561 & 1.113605 & 1.190280 & 1.267813 & 1.385945 \\
OFF\#2 & 1.034427 & 1.068469 & 1.102132 & 1.135420 & 1.200904 & 1.327729 & 1.449538 & 1.624253 \\
OFF\#3 & 1.067034 & 1.132039 & 1.195056 & 1.256139 & 1.372769 & 1.585968 & 1.776429 & 2.028926 \\
OFF\#4 & --- & --- & --- & --- & --- & --- & --- & 1.506337 \\
\bottomrule
\end{tabular}
\end{table}


\section{$c_{1}$의 거울쌍 분해}
\label{app:mirror}

임계선 외 영점 $\rho=\sigma_{0}+it_{0}$에 대해
거울쌍은 $\tilde\rho=1-\sigma_{0}+it_{0}$이다.
거울쌍의 $A'(\rho)/A(\rho)$에 대한 주도 기여는
\[
  \frac{1}{\rho-\tilde\rho} = \frac{1}{2\sigma_{0}-1} = \frac{1}{2d}
\]
으로, $c_{1}^{\text{거울}} = 2\,\mathrm{Re}(1/(2d)) = 1/d$를 준다.
표~\ref{tab:mirror}는 결과~\#117의 분해를 보인다.

\begin{table}[ht]
\centering
\caption{$c_{1}$의 거울쌍 대 배경 분해 (결과~\#117).}
\label{tab:mirror}
\small
\begin{tabular}{lrrrrrr}
\toprule
영점 & $d$ & $c_{1}^{\text{피팅}}$ & $c_{1}^{\text{거울}}=1/d$ & $c_{1}^{\text{배경}}$ & $c_{1}^{\text{배경}}/c_{1}^{\text{거울}}$ \\
\midrule
OFF\#1 & 0.3085 & 3.760 & 3.241 & 2.139 & 1.320 \\
OFF\#2 & 0.1508 & 6.916 & 6.630 & 3.601 & 1.086 \\
OFF\#3 & 0.0744 & 13.543 & 13.449 & 6.818 & 1.014 \\
OFF\#4 & 0.2243 & 5.026 & 4.459 & 2.796 & 1.254 \\
\bottomrule
\end{tabular}
\end{table}

$d\to 0$에 따라 배경 대 거울 비율이 $1$에 접근하여,
두 기여가 극한에서 동일한 크기를 가짐을 나타낸다.
데이터는 $c_{1}\cdot d\to 1$을 보이는데, 이는
배경 항들 사이에 유의미한 상쇄가 있음을 시사한다.


% ============================================================================
\begin{thebibliography}{99}

\bibitem{kim2026}
Kim, Y.-H. (2026).
\textit{Phase Quantisation from Optical Lattices to the Zeta Landscape:
A Unified Framework}.
Preprint, arXiv:2026.xxxxx.

\bibitem{titchmarsh1986}
Titchmarsh, E.~C. (1986).
\textit{The Theory of the Riemann Zeta-Function} (2nd ed., revised by
D.~R.~Heath-Brown).
Oxford University Press.

\bibitem{edwards1974}
Edwards, H.~M. (1974).
\textit{Riemann's Zeta Function}.
Academic Press.

\bibitem{davenport1936}
Davenport, H., \& Heilbronn, H. (1936).
On the zeros of certain Dirichlet series.
\textit{Journal of the London Mathematical Society}, 11, 181--185.

\bibitem{silverman1992}
Silverman, J.~H. (1992).
\textit{The Arithmetic of Elliptic Curves} (2nd ed.).
Springer.

\bibitem{iwaniec2004}
Iwaniec, H., \& Kowalski, E. (2004).
\textit{Analytic Number Theory}.
American Mathematical Society.

\bibitem{epstein1903}
Epstein, P. (1903).
Zur Theorie allgemeiner Zetafunktionen.
\textit{Mathematische Annalen}, 56, 615--644.

\bibitem{rubinstein2001}
Rubinstein, M. (2001).
Numerical computations concerning the Riemann hypothesis.
\textit{PhD thesis}, Princeton University.

\bibitem{cremona2019}
Cremona, J.~E. (2019).
\textit{Algorithms for Modular Elliptic Curves} (2nd ed.).
Cambridge University Press.

\bibitem{pari2023}
PARI/GP (2023). Version 2.15.4.
\url{https://pari.math.u-bordeaux.fr/}.

\bibitem{mpmath2023}
Johansson, F., et al.\ (2023).
\textit{mpmath: a Python library for arbitrary-precision
floating-point arithmetic} (version 1.3.0).
\url{https://mpmath.org/}.

\end{thebibliography}

\end{document}
